% Référence de conception et de modélisation du radar -- outil de préparation de mission C.44
% Traduction française de radar_design.tex (même contenu, même structure)
% Compilation : tectonic radar_design_fr.tex   (ou pdflatex/lualatex deux fois)
\documentclass[11pt,a4paper]{article}

\usepackage[margin=2.3cm]{geometry}
\usepackage[T1]{fontenc}
\usepackage[utf8]{inputenc}
\usepackage{lmodern}
\usepackage[french]{babel}
\usepackage{microtype}
\usepackage{amsmath,amssymb}
\usepackage{icomma}
\usepackage{booktabs}
\usepackage{tabularx}
\usepackage{longtable}
\usepackage{array}
\usepackage{siunitx}
\usepackage{xcolor}
\usepackage{tikz}
\usetikzlibrary{arrows.meta,positioning,shapes.geometric,calc,fit}
\usepackage{enumitem}
\usepackage{caption}
\usepackage[hidelinks]{hyperref}
\usepackage{fancyhdr}

\sisetup{locale=FR,group-separator={\,},group-minimum-digits=4,detect-all}
\setlist{itemsep=2pt,topsep=4pt}
\newcolumntype{L}{>{\raggedright\arraybackslash}X}
\renewcommand{\arraystretch}{1.15}

\newcommand{\code}[1]{\texttt{\small #1}}
\newcommand{\dB}{\,\mathrm{dB}}
\newcommand{\Tsys}{T_{\mathrm{sys}}}
\newcommand{\SNR}{\mathrm{SNR}}
\newcommand{\Pd}{P_{\mathrm d}}
\newcommand{\Pfa}{P_{\mathrm{fa}}}
\newcommand{\Denv}{D_{\mathrm{env}}}

\pagestyle{fancy}
\fancyhf{}
\lhead{\small Radar C.44 -- référence de conception et de modélisation}
\rhead{\small Modèle V6.27.0}
\cfoot{\thepage}

\title{\textbf{Radar AESA spatial à quatre faces}\\[4pt]
\large Référence de conception et de modélisation de l'outil de préparation de mission}
\author{Projet C.44}
\date{27 octobre 2026 -- modèle radar V6.27.0 (étape 5.3 de la feuille de route fermée : P1 verdict, P2 conseil, P3 redondances)}

\begin{document}
\maketitle

\noindent\fbox{\parbox{\dimexpr\linewidth-2\fboxsep-2\fboxrule}{\small\textbf{Conventions.}
(1) Les missions de test du dépôt et leurs sorties figées sont des cas de non-régression.
(2) Les valeurs technologiques (TRM, récepteurs, convertisseurs, oscillateur) sont des hypothèses optimistes pour
2050--2060 : d'une plage crédible, le bout optimiste est retenu. Les valeurs de menace et d'environnement sont des
moyennes de classe ou des médianes.
(3) Toutes les valeurs sont dans \code{radar\_system.toml} ; chaque entrée porte sa valeur, son unité, son statut et sa
raison. Si ce document et le fichier diffèrent, le fichier fait foi.}}

\tableofcontents
\newpage

%=====================================================================================================================
\section{L'outil}
%=====================================================================================================================

\subsection{Ce que fait l'outil}
L'outil est un outil de \emph{préparation de mission}. L'opérateur décrit ce que le radar doit faire pendant une
mission :
\begin{itemize}
  \item des \textbf{volumes de veille} (« secteurs ») : une boîte ou une calotte angulaire, une portée de détection
        exigée, une période de revisite, une menace (classe de cible, surface équivalente radar, vitesse), des réglages
        de résolution et de fausses alarmes, des tranches de distance de poursuite active facultatives ;
  \item le \textbf{ciel} : les corps devant les faces (Soleil, Terre, Lune, \dots), un repère galactique, des
        brouilleurs à distance ;
  \item la \textbf{navigation} : l'attitude du vaisseau et, en option, son orbite ;
  \item des réglages globaux : plafond de puissance facultatif par face (les faces se partagent les \SI{8}{MW} du
        vaisseau), spectre interdit au radar, fenêtre d'élévation.
\end{itemize}
Pour chaque volume, l'outil choisit la fréquence et les formes d'onde de veille en bande S, calcule les portées
obtenues et l'énergie et le temps que chaque face y consacre, ajoute la confirmation en bande Ku de chaque détection et
la poursuite active, et signale ce que le système ne peut pas faire (puissance ou temps sursouscrits, portée non
atteinte, directions aveuglées, limites de dépointage, brouillage d'autoprotection, \dots). La mission est validée
contre toutes les ressources du système radar.

\subsection{Place dans une simulation plus large}
L'outil est la couche \emph{calculateur radar} d'une architecture à quatre couches (figure~\ref{fig:layers}) qui
sépare la \emph{vérité}, connue de la seule simulation, de la \emph{connaissance}, ce que sait le vaisseau :
\begin{enumerate}
  \item \textbf{simulation globale de scénario} : orbites, objets, propagation newtonienne (la vérité) ;
  \item \textbf{convertisseur} : traduit la vérité dans le repère du vaisseau pour la simulation RF, et en message de
        navigation pour le calculateur radar ;
  \item \textbf{calculateur radar} (cet outil) : ne connaît que sa navigation et ses propres détections ; garde les
        volumes attachés au vaisseau, ordonnance, choisit les formes d'onde ;
  \item \textbf{simulation RF} : joue les ordres du calculateur radar contre la vérité.
\end{enumerate}
En préparation de mission, la navigation est un instantané saisi par l'opérateur ; plus tard, ce sera le message de
navigation du simulateur global. Volumes et ciel sont donnés dans un \emph{repère de mission} (vaisseau, inertiel, RTN,
limbe, deux directions) ; le calculateur radar applique la rotation par l'attitude (section~\ref{sec:frames}).

\begin{figure}[ht]
\centering
\begin{tikzpicture}[node distance=6mm and 10mm, box/.style={draw,rounded corners,align=center,minimum width=33mm,
  minimum height=11mm,font=\small}, >={Stealth[length=2mm]}]
\node[box,fill=gray!12] (sim) {1. Simulation globale\\de scénario (vérité)};
\node[box,right=of sim] (conv) {2. Convertisseur};
\node[box,right=of conv,fill=blue!8] (rc) {3. Calculateur radar\\(cet outil)};
\node[box,below=12mm of conv] (rf) {4. Simulation RF};
\draw[->] (sim) -- (conv);
\draw[->] (conv) -- node[above,font=\scriptsize]{navigation} (rc);
\draw[->] (conv) -- node[left,font=\scriptsize,align=right]{vérité dans le\\repère vaisseau} (rf);
\draw[->] (rc) |- node[pos=0.25,right,font=\scriptsize]{ordres} (rf);
\draw[->] (rf.west) -| node[pos=0.75,left,font=\scriptsize]{échos} ($(sim.south)+(0,0)$);
\end{tikzpicture}
\caption{Architecture à quatre couches : vérité et connaissance sont séparées. L'outil réalise la couche 3 en
préparation de mission.}
\label{fig:layers}
\end{figure}

\subsection{Méthode de conception}
Le modèle a été reconstruit pas à pas le long d'une feuille de route (section~\ref{sec:roadmap}). Quelques règles
en façonnent chaque partie :
\begin{itemize}
  \item \textbf{Un seul fichier de paramètres.} Chaque valeur du système est une entrée de \code{radar\_system.toml}
        avec \code{value}, \code{unit}, \code{status} (V : verrouillée par le concepteur ; D : dérivée ou proposée, à
        confirmer), \code{source} (la raison) et \code{date}. Les réglages de l'outil (grilles de recherche, caches,
        valeurs par défaut) sont dans un bloc \code{[tool.*]} séparé. 290 entrées système existent aujourd'hui, dont
        247 verrouillées.
  \item \textbf{Dérivé, pas saisi.} Tout ce qui découle d'autres valeurs est calculé : la composition de l'antenne, les
        diagrammes, la température de bruit du système, le facteur de bruit du convertisseur, l'impulsion la plus
        longue, la grille de fréquences, \dots{} Le chargeur refuse les combinaisons incohérentes avec un message clair.
  \item \textbf{Technologie optimiste, menaces médianes.} Pour notre technologie, le bout optimiste d'une plage
        crédible 2050--2060 ; pour les menaces et l'environnement, la médiane ou la moyenne de classe.
  \item \textbf{Approximations sur décision seulement.} Les raccourcis exacts (formes closes, caches, élagage exact)
        passent d'abord ; depuis le 25 octobre 2026, les approximations sont admises quand elles ne dénaturent pas le
        modèle, chacune décidée par le concepteur avec son effet mesuré et gardée comme un réglage dont la valeur exacte
        reste disponible.
\end{itemize}

%=====================================================================================================================
\section{Le radar}
%=====================================================================================================================

\subsection{Faces et géométrie}
Le radar a \textbf{quatre faces planes} réparties à \SI{90}{\degree} autour de l'axe du vaisseau (azimuts
\SIlist{0;90;180;270}{\degree}), toutes \textbf{inclinées de \SI{25}{\degree} vers le nez}. Le repère du vaisseau a
$Z$ le long de l'axe, vers le nez (modèle « phare » : l'axe est vertical, le nez au zénith) ; l'azimut $\varphi$ est
compté autour de $Z$ de $+X$ vers $+Y$, l'élévation $\varepsilon$ du plan $XY$ vers $+Z$.

Chaque face couvre un \emph{secteur} de rayon \SI{75}{\degree} autour de sa normale (la plus grande limite de
dépointage des bandes), coupé par les plans $\varphi = \pm\SI{45}{\degree}$ entre faces voisines. Les positions de
veille sont comptées en cosinus directeurs $(u, v)$ de la face : un réseau plan dépointe un faisceau de largeur
constante en $(u,v)$, qui s'élargit en $1/\cos\theta$ en angle. Le secteur couvre une aire de $2,0006$ en $(u,v)$.

La géométrie du vaisseau vient d'un modèle 3D à faible nombre de polygones (\SI{149.46}{m} de long, \SI{74.68}{m}
d'envergure sur ses ailes de radiateurs, coque de \SI{17.22}{m} de diamètre) ; les quatre panneaux radar sont des
octogones posés sur un cône autour du poste avant, avec les centres d'ouverture à \SI{5.248}{m} de l'axe. Le chargeur
vérifie le modèle contre les dimensions verrouillées et vérifie que chaque centre d'ouverture est sur un panneau dont
la normale est celle de la face à l'inclinaison verrouillée.

\subsection{Bandes et rôles}
\begin{center}
\begin{tabularx}{\linewidth}{@{}lLL@{}}
\toprule
 & \textbf{Bande S} & \textbf{Bande Ku}\\
\midrule
Bande utilisable & \SIrange{2.60}{2.93}{GHz} & \SIrange{13.5}{14.5}{GHz}\\
Bande c\oe{}ur (sans perte de bord) & \SIrange{2.66}{2.88}{GHz} & \SIrange{13.7}{14.3}{GHz}\\
Rôle & veille (détection de nouvelles cibles) & confirmation des détections, poursuite active\\
Limite de dépointage & \SI{75}{\degree} jusqu'à \SI{2.85}{GHz}, \SI{70}{\degree} au-delà & \SI{75}{\degree} jusqu'à \SI{14.25}{GHz}, \SI{70}{\degree} au-delà\\
TRM par face & \num{4900} & \num{122500}\\
Puissance crête par TRM (sortie PA) & \SI{2000}{W} & \SI{75}{W}\\
Puissance crête par face & \SI{9.80}{MW} & \SI{9.19}{MW}\\
Bande instantanée & \SI{200}{MHz} & \SI{200}{MHz}\\
Température système, ciel froid & \SI{144.5}{K} & \SI{230.8}{K}\\
\bottomrule
\end{tabularx}
\end{center}

L'architecture est \textbf{entièrement numérique} : un convertisseur analogique-numérique par TRM et par voie de
polarisation, retard vrai numérique dans chaque TRM (pas de dépointage en fréquence, pas de limite de remplissage
d'ouverture sur la bande), faisceaux de réception formés par le calcul. La veille utilise la seule bande S ; le
« toucher » S d'un mode de confirmation est l'unique exception.

\subsection{Composition de l'antenne}
Chaque face est un pavage bibande de \textbf{196 TRSA} (sous-réseaux d'émission-réception). Une TRSA est une tuile
carrée de $25\times25$ TRM Ku (pas \SI{10.4}{mm}) et $5\times5$ TRM S (pas \SI{52}{mm}, proche de $\lambda/2$ en S) ;
chaque TRM S couvre exactement un bloc de $5\times5$ TRM Ku. Les TRSA sont disposées en octogone approché de 16 rangées
centrées,
\[
  [6, 8, 10, 12, 14, 16, 16, 16, 16, 16, 16, 14, 12, 10, 8, 6],
\]
séparées par un jeu mécanique de \SI{0.5}{mm} (zone morte, non rayonnante). Grandeurs dérivées :
\begin{center}
\begin{tabularx}{\linewidth}{@{}Lr@{}}
\toprule
Côté d'une TRSA & \SI{260}{mm}\\
Surface émissive & \SI{13.2496}{m^2} (remplissage \SI{99.65}{\percent} de \SI{13.2964}{m^2})\\
Enveloppe $\Denv$ (rangée la plus large : 16 TRSA et 15 jeux) & \SI{4.1675}{m}\\
Pas moyen des éléments (pas des TRSA / TRM par côté) & S \SI{52.1}{mm}, Ku \SI{10.42}{mm}\\
Étendue RMS de l'ouverture selon un axe $x_{\mathrm{rms}}$ & \SI{1.033}{m} $= 0,248\,\Denv$\\
\bottomrule
\end{tabularx}
\end{center}
Le jeu a été ramené de \SI{2}{mm} à \SI{0.5}{mm} parce que la rupture du pas des éléments d'une TRSA à la suivante crée
des « lobes de jeu » près du lobe de réseau (\SI{-16.6}{dB} en Ku avec \SI{2}{mm}) qu'aucune pondération de réception ne
peut abaisser ; \SI{0.5}{mm} suppose des matériaux de TRSA et de plaque arrière appariés en dilatation thermique.

\begin{figure}[ht]
\centering
\begin{tikzpicture}[scale=0.32]
  \foreach \n [count=\r from 0] in {6,8,10,12,14,16,16,16,16,16,16,14,12,10,8,6} {
    \pgfmathsetmacro{\xo}{-\n/2}
    \foreach \k in {1,...,\n} {
      \pgfmathsetmacro{\x}{\xo+\k-1}
      \draw[fill=blue!10,draw=blue!50!black,line width=0.2pt] (\x,7.5-\r) rectangle ++(0.94,0.94);
    }
  }
  \draw[->,thick] (0.47,8.4) -- (0.47,10.2) node[above,font=\scriptsize]{$v$ ($e_v$, vers le plan de l'axe du vaisseau)};
  \draw[->,thick] (8.8,0.47) -- (10.8,0.47) node[right,font=\scriptsize]{$u$};
  \draw[|<->|] (-8,-9.2) -- (8,-9.2) node[midway,below,font=\scriptsize]{$\Denv = \SI{4.1675}{m}$ (16 TRSA)};
\end{tikzpicture}
\caption{Les 196 TRSA d'une face (chaque carré : $25\times25$ TRM Ku et $5\times5$ TRM S).}
\label{fig:octagon}
\end{figure}

%=====================================================================================================================
\section{Modèle d'antenne}
%=====================================================================================================================

\subsection{Diagramme calculé à partir de la composition}
Le diagramme est calculé à partir des positions des éléments de la composition (\code{core/antenna.py}). Les éléments
sont isotropes dans le facteur de réseau ; le diagramme élémentaire n'entre que par l'aire projetée $\cos\theta$ et la
désadaptation active (section~\ref{sec:mismatch}). Près de l'axe, en unités de $\lambda/\Denv$, le diagramme ne dépend
pas de la fréquence : un calcul par bande sert toute la bande (S et Ku échantillonnent le même octogone et concordent à
\SI{0.1}{\percent}). Les rangées sont centrées et palindromiques, donc le facteur de réseau est réel,
\[
  AF(u,v) = \sum_{\text{quadrant}} w\,m\,\cos(k x u)\cos(k y v),
\]
avec des multiplicités $m$ pour la symétrie ; avec la symétrie $x\leftrightarrow y$ de l'octogone, les coupes de
\SIrange{0}{45}{\degree} suffisent, chacune un produit matriciel. Les résultats sont mis en cache sous une empreinte de
leurs entrées.

\begin{itemize}
  \item \textbf{Faisceau d'émission} : éclairement uniforme (PA saturés). Largeur à mi-puissance (plus petite coupe)
        $1,036\,\lambda/\Denv$, soit \SI{1.525}{\degree} à \SI{2.8}{GHz} et \SI{0.305}{\degree} à \SI{14}{GHz} ;
        premier lobe secondaire \SI{-15.6}{dB}.
  \item \textbf{Faisceau de réception} : loi de Taylor circulaire (Hansen), $\bar n = 5$, dont le niveau de conception
        est \emph{ajusté numériquement} (fausse position) pour que le niveau \emph{obtenu} sur l'octogone réel soit
        l'option choisie pour le volume de veille : \SIlist{-25;-30;-35}{dB}. Rendement et largeur en découlent :
        \begin{center}
        \begin{tabular}{@{}lccc@{}}
        \toprule
        Option & niveau de conception & $\eta_{\mathrm{RX}}$ & largeur ($\lambda/\Denv$)\\
        \midrule
        \SI{-25}{dB} (défaut) & \SI{29.5}{dB} & 0,873 & 1,173\\
        \SI{-30}{dB} & \SI{35.7}{dB} & 0,773 & 1,255\\
        \SI{-35}{dB} & \SI{41.7}{dB} & 0,693 & 1,325\\
        \bottomrule
        \end{tabular}
        \end{center}
        Baisser le niveau rejette les sources fortes hors de l'axe (cible brillante, Lune ou Terre, brouilleur) au prix
        du gain de réception ($-0,53$ et $-1,00\dB$ par rapport à l'option \SI{-25}{dB}). Un réglage par volume sert S
        et Ku.
  \item \textbf{Gain} : $G = 4\pi A\,\eta\,\eta_{\mathrm{ohm}}\cos\theta/\lambda^2$ avec la surface émissive $A$ ; à
        \SI{2.8}{GHz} dans l'axe, \SI{41.62}{dBi} en émission et \SI{41.03}{dBi} en réception (\SI{-25}{dB}) ; gain
        d'émission Ku à \SI{14}{GHz} : \SI{55.60}{dBi}.
  \item \textbf{Précision angulaire} : la borne de Cramér--Rao de l'ouverture numérique complète avec une marge de
        1,13,
        \[ \sigma_u = 1,13\,\frac{\lambda}{2\pi\,x_{\mathrm{rms}}\sqrt{2\,\SNR}}, \]
        indépendante de la pondération de réception (l'estimateur du maximum de vraisemblance utilise toute
        l'ouverture) ; la même marge s'applique aux estimateurs de Doppler et de distance.
\end{itemize}

\subsection{Lobes lointains et lobes de jeu}
Au-delà de la fenêtre proche ($8\,\lambda/\Denv$), le facteur de réseau est balayé sur une grille 2D de toute la région
visible de la bande (rayon $(1+\sin\theta_{\lim})\Denv/\lambda$ au haut de la bande), voisinages des lobes de réseau
exclus. Un profil cumulé garde le niveau du plus fort lobe à moins de $r$ de l'axe du faisceau ; un lobe à $r$ devient
visible dès que le faisceau est dépointé à $\theta$ avec $r\lambda/\Denv \le 1+\sin\theta$ (pire azimut). Les plus
forts lobes lointains, lobes de jeu compris, sont à \SI{-25.3}{dB} en Ku et \SI{-30.7}{dB} en S (uniforme) ; la
pondération n'abaisse pas les lobes de jeu. Pour chaque secteur, l'outil indique jusqu'à quel dépointage l'option de
réception est tenue sur la région visible ; en Ku, les options \SIlist{-30;-35}{dB} ne sont tenues que jusqu'à
\SIlist{75;62}{\degree} à \SI{14}{GHz}.

Le même module fournit les \emph{moyennes azimutales} du diagramme de puissance de réception en fonction de la
distance à l'axe du faisceau (fines près de l'axe, moyennes par anneaux au-delà). Elles servent à l'intégrale du bruit
du ciel, au fouillis par les lobes secondaires et à la captation des brouilleurs (sections~\ref{sec:sky}
et~\ref{sec:env}). Par le théorème de Parseval, la normalisation est exacte : un corps remplissant tout le demi-espace
avant ajoute sa température de brillance (vérifié à 1,00--1,04).

\subsection{Désadaptation active et lobes de réseau}\label{sec:mismatch}
Un modèle, appliqué à S et Ku, en émission et en réception, décrit la réflexion active du réseau : une nappe de
courant adaptée à $\theta_m = \SI{45}{\degree}$ par une couche d'adaptation grand angle (WAIM), et la montée de la
réflexion à l'approche du premier lobe de réseau dans la région visible ($g_0 = 0,02$, $d_{\mathrm{mean}}$ le pas moyen
des éléments, pire azimut) :
\begin{align*}
  |\Gamma_{\mathrm{WAIM}}| &= \frac{|\cos\theta - \cos\theta_m|}{\cos\theta + \cos\theta_m},\qquad
  |\Gamma_{\mathrm{GL}}| = \exp(-g/g_0),\quad g = \frac{\lambda}{d_{\mathrm{mean}}} - \sin\theta - 1,\\
  1 - |\Gamma|^2 &= (1 - |\Gamma_{\mathrm{WAIM}}|^2)(1 - |\Gamma_{\mathrm{GL}}|^2),
\end{align*}
avec une perte par passage $-10\log_{10}(1-|\Gamma|^2)$, comptée deux fois (émission et réception). À \SI{2.8}{GHz},
l'aller-retour coûte \SI{0.26}{dB} dans l'axe, 0 à \SI{45}{\degree}, \SI{1.12}{dB} à \SI{70}{\degree} et
\SI{2.11}{dB} à \SI{75}{\degree}. Le haut de la bande S est fixé par ce mur : à \SI{2.93}{GHz} et \SI{70}{\degree}, la
perte est de \SI{0.97}{dB} par passage ; un lobe de réseau devient visible à \SI{3.0}{GHz}. L'optimiseur signale toute
tranche utilisée au-delà de \SI{3}{dB} par passage.

La polarisation d'émission est linéaire selon $e_v$ (l'axe de la face dans le plan de l'axe du vaisseau), un PA par
TRM ; sa pureté (polarisation croisée de Ludwig-3, $(1+\cos\theta)/(1-\cos\theta)$ dans le plan à \SI{45}{\degree}) est
donnée en sortie. Avec le récepteur bipolarisation, elle ne coûte rien au bilan.

\subsection{Sous-ouvertures de la partition d'émission}\label{sec:subap}
Pour le mode multifaisceau par partition de l'ouverture (section~\ref{sec:multibeam}), la face est coupée en moitiés
(coupe horizontale \code{h} : rangées au-dessus du centre ; coupe verticale \code{v} : colonnes à droite du centre) ou
en quadrants, chacun uniforme. Leurs diagrammes 2D sont évalués directement (une moitié n'est pas centrée, aucune
symétrie n'est supposée), exactement, par une somme séparable par rangées d'éléments :
\begin{center}
\begin{tabular}{@{}lccc@{}}
\toprule
Ouverture & TRM S & largeurs à $-3\dB$ $u\times v$ ($\lambda/\Denv$) & maille / face complète\\
\midrule
Face complète & \num{4900} & $1,041\times1,041$ & 1\\
Moitié & \num{2450} & $1,041\times1,928$ & $\times1,85$\\
Quadrant & \num{1225} & $1,928\times1,928$ & $\times3,43$\\
\bottomrule
\end{tabular}
\end{center}
Le rapport des mailles est inférieur à 2 et 4 parce qu'un demi-octogone n'est pas deux fois moins large en largeur
effective. L'octogone étant symétrique en $x$ et $y$, les deux moitiés donnent le même faisceau tourné de
\SI{90}{\degree} : le code le vérifie exactement et n'en cherche alors qu'une.

%=====================================================================================================================
\section{Chaîne d'émission et bilan de puissance}
%=====================================================================================================================

\subsection{Modèle de TRM}
Les TRM sont en GaN sur diamant avec refroidissement par microcanaux. La puissance crête est spécifiée à la sortie du
PA (où le rendement en puissance ajoutée, PAE, est défini) ; le chemin du PA à l'élément est une perte dérivée (filtre
d'harmoniques, circulateur, coupleur d'étalonnage, ligne : \SI{0.6}{dB} en S, \SI{0.95}{dB} en Ku).

\paragraph{PAE.} Pour une impulsion de durée $\tau$ dans un train de facteur de charge $d = \tau/\mathrm{PRI}$,
\[
  \eta = \eta_{\mathrm{sat}}\,\frac{\tau}{\tau + t_{\mathrm{edge}}}\,\min\bigl(1,\;1 - a\,(d - d_{\mathrm{ref}})\bigr),
\]
avec la PAE saturée $\eta_{\mathrm{sat}}$ au facteur de charge de référence $d_{\mathrm{ref}} = 0,10$, le temps
$t_{\mathrm{edge}}$ pendant lequel le PA consomme sans RF pleine (commutation, établissement de la polarisation et de
la phase), et une pente $a$ pour l'échauffement du canal. Une impulsion seule utilise $d_{\mathrm{ref}}$.

\begin{center}
\begin{tabular}{@{}lcc@{}}
\toprule
 & S & Ku\\
\midrule
$\eta_{\mathrm{sat}}$ & 0,80 & 0,65\\
$t_{\mathrm{edge}}$ & \SI{0.2}{\micro s} & \SI{0.1}{\micro s}\\
Pente $a$ & 0,05 ($-2,75\,\%$ à $d=0,65$) & 0,08 ($-4,4\,\%$ à $d=0,65$)\\
Plafond du facteur de charge des trains & 0,65 & 0,65\\
Réserve d'énergie par TRM & \SI{25}{J} & \SI{1.2}{J}\\
Impulsion la plus longue $\tau_{\max}$ & \SI{2.0}{ms} & \SI{2.08}{ms}\\
\bottomrule
\end{tabular}
\end{center}

\paragraph{Réserve d'énergie.} Chaque TRM tire son impulsion d'une batterie de condensateurs locale. Une impulsion qui
en utilise une fraction $x = E/E_{\mathrm{buf}}$ perd en moyenne $-10\log_{10}(1-x/2)$ (la sortie du PA suit l'énergie
stockée) ; une impulsion au-delà de $x = 0,2$ est refusée, ce qui fixe
$\tau_{\max} = 0,2\,E_{\mathrm{buf}}\,\eta_{\mathrm{sat}}/P_{\mathrm{pk}} - t_{\mathrm{edge}}$. L'impulsion la plus longue
finit \SI{0.97}{dB} plus bas (moyenne \SI{0.46}{dB}) ; la conversion AM/PM du PA pendant que la réserve se vide est
donnée en sortie, non comptée (dérive de phase quasi linéaire).

\subsection{Bilan électrique}\label{sec:power}
Le bilan d'une face est \emph{électrique} : chaque visée coûte son énergie RF divisée par la PAE de son impulsion, et
la chaleur de la face est l'énergie électrique moins l'énergie RF (environ \SI{370}{W} par TRM S pendant un long train
à $d = 0,65$, une entrée de l'étape thermique). Depuis l'étape 4.1, la puissance électrique est modélisée à quatre
niveaux.

\paragraph{Niveau vaisseau : un budget partagé.} Le vaisseau fournit \SI{8}{MW} au radar. Les faces n'ont \emph{pas}
de parts fixes : elles puisent ensemble dans ce budget, et une face très chargée peut prendre plus d'un quart quand les
autres en demandent moins. Le verdict d'énergie est donc porté au niveau du vaisseau,
\[
  \sum_{k=1}^{4} P_{\mathrm{bus},k} \le \SI{8}{MW},\qquad
  P_{\mathrm{bus},k} = \frac{P_{\mathrm{TX},k} + P_{\mathrm{on}} + p_{\mathrm{pump}}\,Q_k}{\eta_{\mathrm{psu}}\,\eta_{\mathrm{dist}}},
\]
où $P_{\mathrm{TX},k}$ est la puissance électrique moyenne des émissions de la face $k$ (veille, confirmation,
poursuite, à l'alimentation des TRM), $P_{\mathrm{on}}$ la consommation permanente, $Q_k$ la chaleur moyenne de la face
et $p_{\mathrm{pump}} = 2\,\%$ la part de cette chaleur dépensée par les pompes ; $\eta_{\mathrm{psu}} = 0,95$ et
$\eta_{\mathrm{dist}} = 0,995$ sont les rendements de l'alimentation et de la distribution (produit 0,945). Chaque face
est aussi limitée par ce que peut porter sa propre alimentation (câbles, convertisseur), \SI{8}{MW}, et par un plafond
facultatif que l'opérateur peut lui réserver dans la mission (une valeur pour toutes les faces ou une par face).
L'outil donne la puissance que chaque face tire réellement, la marge du budget partagé et la répartition entre faces
(section~\ref{sec:load}). Pour arbitrer énergie et temps dans le choix des formes d'onde, l'optimiseur de veille garde
une répartition de \emph{référence} de $\SI{8}{MW}/4 = \SI{2}{MW}$ par face.

\paragraph{Consommation permanente.} Un réseau numérique reçoit en permanence, qu'il émette ou non. Par face : deux
convertisseurs par TRM (un par polarisation ; \SI{0.1}{W} chacun à \SI{250}{MS/s} et 13 bits, valeur optimiste 2050,
contre 0,3--1 W aujourd'hui) et deux chaînes de réception par TRM (LNA, amplificateurs, distribution de l'oscillateur
local : \SI{0.5}{W} en S, \SI{0.1}{W} en Ku), dans les deux bandes ; \SI{50}{kW} de formation de faisceaux et de
traitement numériques (plafond, section~\ref{sec:compute}) ; \SI{2}{kW} de contrôle. Total \SI{106.9}{kW} par face
(\num{9800} convertisseurs S et \num{245000} Ku). Avec la référence de \SI{2}{MW}, il reste à une face
$2\times0,945 - 0,107 = \SI{1.784}{MW}$ pour ses émissions.

\paragraph{Niveau TRM : alimentation sur un train.} Chaque TRM tire son impulsion de sa réserve locale et la recharge
par une alimentation continue (\SI{1700}{W} en S, \SI{80}{W} en Ku). Sur un train de $n$ impulsions d'énergie
$E_p = P_{\mathrm{pk}}\tau/\eta$ à chaque PRI, deux conditions doivent tenir : la puissance moyenne $E_p/\mathrm{PRI}$
reste dans l'alimentation, et la réserve, rechargée au débit de l'alimentation entre les impulsions, perd au plus la
fraction admise :
\[
  (n-1)\,\max\bigl(0,\;E_p - P_{\mathrm{sup}}\,\mathrm{PRI}\bigr) + E_p \le 0,2\,E_{\mathrm{buf}}.
\]
Une forme d'onde qui viole l'une des conditions est retirée des espaces de recherche (aucune forme d'onde S
aujourd'hui ; 68 visées Ku sur \num{2356}, trains longs à fort facteur de charge).

\paragraph{Niveau face : stockage et puissance crête.} Quand une face émet, tous ses TRM consomment ensemble : environ
\SI{12.3}{MW} pour une visée S à pleine puissance ($\num{4900}\times\SI{2}{kW}/0,8$) et \SI{14.1}{MW} en Ku, bien
au-dessus de la moyenne que délivre le bus. L'écart vient d'un stockage d'énergie par face (supercondensateurs,
\SI{5}{MJ}, \SI{80}{\percent} utilisables, décharge \SI{15}{MW}), rechargé au débit que laisse l'alimentation après la
consommation permanente. Trois contrôles :
\begin{itemize}
  \item \emph{une observation} : l'énergie de l'observation la plus coûteuse du plan (un passage de veille S ou une
        visée Ku), moins la recharge pendant celle-ci, doit tenir dans le stockage utilisable ;
  \item \emph{pire rafale} : si toutes les émissions de la plus longue revisite de la face étaient groupées, le
        déficit $RR\,(P_{\mathrm{TX}} - P_{\mathrm{rech}}\,f_t)$ ($f_t$ la part du temps en émission) doit tenir dans
        le stockage utilisable ; il ne tient jamais pour des revisites longues, d'où une \emph{règle d'ordonnanceur} :
        au plus $E_{\mathrm{usable}}/(P_{\mathrm{peak}} - P_{\mathrm{rech}})$ d'émission contiguë à puissance crête
        (environ \SI{0.3}{s}), puis recharge ;
  \item \emph{puissance crête} : la crête d'une observation doit rester dans la décharge du stockage plus la recharge.
\end{itemize}
L'échelon de charge du bus au début d'une observation et sa pente sur le front des impulsions (\SI{0.2}{\micro s} en
S) sont donnés en sortie pour la conception du bus.

\subsection{Bilan thermique}\label{sec:thermal}
Toute la puissance électrique tirée par une face devient chaleur, sauf l'énergie RF rayonnée. Depuis l'étape 4.2, la
chaleur est contrôlée à trois échelles, chacune avec sa constante de temps.

\paragraph{Fluide et plaque froide.} Chaque face est refroidie par des microcanaux dans sa plaque froide, avec un
fluide diélectrique (dérivé du submer SmartCoolant ; valeurs 2020 : capacité calorifique \SI{2.26}{kJ/(kg.K)},
coefficient d'échange d'environ \SI{2050}{W/(m^2.K)} en circulation forcée, densité 0,796). Hypothèse 2050--2060 :
\SI{2.5}{kJ/(kg.K)}, \SI{3000}{W/(m^2.K)}, 0,8. Les microcanaux mouillent dix fois l'empreinte ; la plaque peut être
\SI{40}{K} au-dessus d'un fluide à \SI{20}{\celsius}. D'où un flux admissible de
$h \times 10 \times \SI{40}{K} = \SI{120}{W/cm^2}$ et une capacité de refroidissement d'environ \SI{15.9}{MW} par face :
bien au-delà de ce que dissipe une face (un TRM en long train S dégage environ \SI{14}{W/cm^2} sur son empreinte ; une
face qui tire les \SI{8}{MW} dissipe environ \SI{2}{MW}). Toutes ces valeurs sont dérivées du fichier : un changement
de spécification du fluide est absorbé.

\paragraph{Chaleur d'une face.} Le circuit d'une face porte la chaleur des émissions (énergie électrique moins RF),
de l'électronique permanente (\SI{106.9}{kW}) et des pompes ; les pertes de conversion de puissance vont au vaisseau.
Une rafale d'émissions est absorbée par l'inertie thermique de la face (\SI{50}{L} de fluide et environ \SI{100}{kg} de
structure, \SI{190}{kJ/K} ; \SI{10}{K} de montée admise, un tampon de \SI{1.9}{MJ}) ; l'outil contrôle la pire rafale
et donne l'émission contiguë la plus longue à chaleur crête comme règle d'ordonnanceur (aucune n'est nécessaire
aujourd'hui).

\paragraph{Jonction des amplificateurs.} La température de jonction d'un amplificateur de TRM suit un réseau de Foster
à deux pôles (puce et diffuseur diamant : S \SI{0.02}{K/W} à \SI{10}{\micro s} et \SI{0.03}{K/W} à \SI{1}{ms} ; Ku
\SI{0.2}{K/W} et \SI{0.3}{K/W}) excité par la chaleur de chaque impulsion, plus la convection vers le fluide à la
chaleur moyenne. Sur toutes les formes d'onde de veille et les visées Ku, elle reste sous \SI{63}{\celsius}, loin de
la limite de \SI{225}{\celsius}, même en émission continue : l'échauffement ne limite pas le facteur de charge, dont le
plafond de 0,65 traduit les limites de rendement et d'alimentation. La température est donnée par forme d'onde ; le
retrait des formes d'onde au-delà de la limite existe comme réglage et est désactivé.

\paragraph{Vaisseau.} La chaleur des quatre faces plus les pertes de conversion doit tenir dans le rejet thermique que
le vaisseau réserve au radar, partagé par les faces comme le budget électrique (\SI{6}{MW}, provisoire).

%=====================================================================================================================
\section{Modèle de bande}
%=====================================================================================================================
Chaque bande a une bande utilisable et une bande c\oe{}ur. Un signal de largeur $B$ centré sur $f$ occupe
$[f-B/2,\,f+B/2]$, qui doit tenir dans la bande utilisable (et hors du spectre interdit par la mission) ; sinon la
forme d'onde n'est pas jouable à $f$. Hors de la bande c\oe{}ur, une dégradation en cosinus surélevé atteint
\texttt{edge\_derating\_db} aux bords utilisables (\SI{1.0}{dB} aller-retour en S, \SI{0.5}{dB} en Ku) ; sa moyenne
\emph{exacte} sur le spectre occupé est appliquée. Le chargeur vérifie que la bande c\oe{}ur et la fréquence de rupture de
la limite de dépointage sont dans la bande utilisable : la bande se change dans le seul fichier. Avec le retard vrai
numérique, il n'y a pas de limite de remplissage d'ouverture (un dépointage par la phase limiterait le signal à
$c/(\Denv\sin\theta)$, environ \SI{75}{MHz} à \SI{75}{\degree}).

%=====================================================================================================================
\section{Chaîne de réception}\label{sec:rx}
%=====================================================================================================================

\subsection{Température de bruit du système}
$\Tsys$ est la cascade de Friis ramenée à l'accès de l'antenne : température d'antenne en ciel froid, pertes avant le
LNA (circulateur, limiteur, ligne : elles atténuent et ajoutent $(L-1)T_{\mathrm{phys}}$), le LNA, la chaîne après le
LNA et le CAN, chacun divisé par le gain qui le précède, et un facteur de bruit numérique. Valeurs (S / Ku) : ciel froid
\SI{3.2}{K} / \SI{2.8}{K}, tête de réception à \SI{320}{K}, pertes avant LNA \SI{0.85}{dB} / \SI{1.2}{dB}, facteur de
bruit du LNA \SI{0.75}{dB} / \SI{1.15}{dB}, soit \textbf{\SI{144.5}{K} en S et \SI{230.8}{K} en Ku} en ciel froid. Le
ciel devant la face est ajouté par direction (section~\ref{sec:sky}).

\subsection{Convertisseur et dynamique}
Le facteur de bruit du CAN est dérivé du convertisseur :
\[
  \mathrm{NF}_{\mathrm{ADC}} = P_{\mathrm{FS}} - (6,02\,\mathrm{ENOB} + 1,76) - 10\log_{10}(k T_0 f_s),
\]
avec 13 bits effectifs, une pleine échelle de \SI{+10}{dBm} et $f_s = 1,25\times\SI{200}{MHz} = \SI{250}{MS/s}$
complexes : \SI{20.0}{dB}. Les niveaux sont ramenés à l'accès de l'antenne : le niveau d'écrêtage est
$P_{\mathrm{FS}} - \SI{3}{dB}$ moins le gain jusqu'au CAN (\SI{53.4}{dB} au-dessus du bruit en S). La numérisation se
fait par élément, \emph{avant} la formation des faisceaux : une source forte dans n'importe quelle direction avant
atteint chaque convertisseur. La \emph{portée minimale sans saturation} est celle où l'écho crête par élément d'une
visée à pleine puissance et un seul faisceau atteint le niveau d'écrêtage (S : \SI{12.2}{km} pour la cible de
référence, \SI{2.3}{km} pour un KKV). Le récepteur est sourd pendant $\tau$ plus la récupération du limiteur
(\SI{50}{ns} en S, \SI{25}{ns} en Ku).

\subsection{Double polarisation}
Chaque TRM a deux voies de réception. Une visée de détection utilise un détecteur GLRT sur la polarisation inconnue de
l'écho, estimée sur la visée : une inconnue complexe de plus, une perte $\SNR_{\mathrm{req}}(n+1)-\SNR_{\mathrm{req}}(n)$
(\SI{0.47}{dB} pour un échantillon cohérent, faible pour les longs trains non cohérents). Une visée de poursuite connaît
la polarisation par la piste (\SI{0.05}{dB} de résidu). Le récepteur bivoie recueille aussi l'écho en polarisation
croisée de la cible, un gain $10\log_{10}(1+\rho)$ avec $\rho$ le rapport des SER croisée et principale, et les deux
voies ajoutent une diversité de polarisation à la fluctuation de la cible (section~\ref{sec:targets}).

\subsection{Isolation entre bandes et entre faces}
La fuite de chaque bande et face en émission vers chaque élément en réception est contrôlée contre trois limites :
blocage du LNA, écrêtage du CAN et désensibilisation dans le canal (\SI{6}{dB} sous le bruit). Avec les couplages
retenus (\SI{90}{dB} vers une face adjacente, \SI{120}{dB} vers la face opposée, \SI{-30}{dB} à l'intérieur d'une
TRSA), une face ne peut pas recevoir une bande pendant qu'une autre face l'émet : toutes les faces partagent
\textbf{une cadence commune d'émission et de réception par bande}. Règle supplémentaire : pas de réception Ku sur une
face pendant une émission S dont la cinquième harmonique $[5(f-B/2),\,5(f+B/2)]$ couvre le canal Ku. Ce sont des
règles de l'ordonnanceur (section~\ref{sec:sched}).

\subsection{Bruit du ciel}\label{sec:sky}
La température d'antenne dépend de la direction visée. Avec le gain de réception écrit $G = G_0\,w\,g(\rho)$
(diagramme élémentaire $w=\cos\theta$, rien derrière la face) et $d\Omega = du\,dv/w$, le facteur $w$ se simplifie et
la montée vaut
\[
  \Delta T = T_b\,\frac{A\,\eta_{\mathrm{RX}}\,\eta_{\mathrm{ohm}}}{\Denv^2}\int_{\text{corps}} g(\rho)\,d\rho^2,
\]
l'intégrale portant sur le corps vu dans le plan $(u,v)$ de la face (moitié avant seulement), $\rho$ étant la distance
à l'axe du faisceau en $\lambda/\Denv$ et $g$ la moyenne azimutale du diagramme de réception calculé, lobe principal
compris. Les petits corps (Soleil, Lune) sont échantillonnés finement sur leur disque ; les grands (Terre) par une
grille polaire dans la fenêtre proche et des anneaux au-delà. Les températures de brillance sont des valeurs de classe
par bande (Soleil calme, mi-cycle et actif : \SIlist{4.5e4;1.0e5;1.5e5}{K} en S ; Terre \SI{290}{K} au pire, variantes
océan et terre ; Lune \SI{230}{K}) ; le fond galactique suit un modèle plan et bulbe en $f^{-2,5}$, avec un pire cas
appliqué partout quand la mission ne donne pas de repère galactique. Le Soleil dans le faisceau ajoute environ
\SI{17}{dB} en S comme en Ku ; la Lune environ \SI{2.6}{dB} en Ku.

Pour un secteur, $\Tsys$ est évaluée tous les demi-lobes d'émission sur les positions balayées (trame de
\SI{0.2}{\degree}). Près des sources chaudes, le bruit varie plus vite que ce pas : les cases de la trame sont évaluées
une à une là où la montée de leur bloc dépasse \SI{0.5}{dB}, et à moins de trois largeurs de faisceau d'émission d'un
brouilleur à distance, du centre d'un petit corps ou du limbe d'un grand corps ; chaque case est alors découpée en
$8\times8$ points de même poids. Un point dont la montée dépasse \SI{10}{dB} est \emph{aveuglé} : la case n'est balayée
que pour la part de ses points restants, et la part aveuglée est donnée avec le corps responsable. La $\Tsys$ de
\emph{budget} du secteur est le quantile $(1-\text{\code{short\_tol}})$ (\SI{99}{\percent}) des positions non aveuglées,
chaque point pour son poids, et toutes ses portées varient en $(T_{\mathrm{cold}}/T_{\mathrm{budget}})^{1/4}$. Tourner
ensemble le ciel et le secteur autour de l'axe du vaisseau d'angles hors de la trame déplace le budget et la part
aveuglée de moins de \SI{1}{\percent} (Soleil, \SI{0.37}{\degree} et \SI{1.13}{\degree}) ; une dérive régulière aux
plus grands angles (\SI{1.4}{\percent} à \SI{2.9}{\degree}), attribuée aux faces fixes, reste à confirmer. La même
règle s'applique en Ku à la confirmation et à la poursuite du secteur.

%=====================================================================================================================
\section{Formes d'onde}\label{sec:wf}
%=====================================================================================================================

\subsection{Modulation et compression d'impulsion}
Deux familles sont modélisées : la modulation linéaire de fréquence (LFM, tolérante au Doppler) et la modulation non
linéaire (NLFM, lobes secondaires en distance bas obtenus par la loi de fréquence, sans perte de pondération). Les codes
de phase sont exclus : faible tolérance au Doppler aux vitesses de rapprochement jusqu'à \SI{25}{km/s}
($f_d \approx \SI{470}{kHz}$ en S). Pour chaque visée, avec $f_{d,\max} = 2v/\lambda$ :
\begin{itemize}
  \item si $f_{d,\max}/B \le 2\,\%$ : NLFM, perte de désadaptation \SI{0.1}{dB} ;
  \item sinon : LFM avec une pondération de Taylor ($\bar n = 5$) en réception seule au niveau de lobes secondaires du
        volume ; perte = perte de pondération $-\,20\log_{10}(1-f_d/B)$ (perte Doppler au sommet de la fonction
        d'ambiguïté).
\end{itemize}
Le niveau des lobes secondaires en distance est l'option angulaire du volume : une politique de lobes secondaires par
volume. Une pondération à \SIlist{-25;-30;-35}{dB} coûte \SIlist{0.41;0.68;0.93}{dB} et élargit le lobe principal de
$\times$1,18, 1,27, 1,34. La résolution physique en distance est cet élargissement fois $c/2B$ (2 m demandent
\SI{89}{MHz} à \SI{-25}{dB}) ; la constante de l'estimateur de distance est $k_r = 0,6/(\text{rapport de largeur de
bande RMS})$. Le biais de couplage distance--Doppler $c f_d \tau/2B$ est corrigé par la mesure Doppler.

\subsection{Effets intra-impulsion et chronologie}
Une cible traverse $n_c = 2v\tau B/c$ cases distance pendant l'impulsion ; un filtre adapté à une autre vitesse a alors
une erreur de phase quadratique, une perte $-20\log_{10}\bigl|\langle e^{\,j2\pi n_c x^2}\rangle_{x\in[-1/2,1/2]}\bigr|$
(1 dB à $n_c = 1,016$). Elle fixe le pas du banc de vitesses (section~\ref{sec:proc}).

Le récepteur est sourd pendant l'impulsion et la récupération du limiteur : distance aveugle
$c(\tau + t_{\mathrm{rec}})/2$ (\SI{3}{km} pour \SI{20}{\micro s}, \SI{300}{km} pour \SI{2}{ms}), qui doit rester sous
la plus courte distance d'intérêt du secteur (\code{range\_min\_km}, 10\,\% de la portée exigée par défaut). Un train
est non ambigu en distance si $\mathrm{PRI} \ge 2R/c + \tau + t_{\mathrm{rec}}$ ; un secteur peut autoriser des trains
ambigus, décalés sur 3 PRI à $\pm10\,\%$, avec la perte moyenne d'éclipse $-10\log_{10}(1-(\tau+t_{\mathrm{rec}})/\mathrm{PRI})$
(ils sont retirés quand une surface est dans le lobe principal, son fouillis se replierait sur les cibles).

\subsection{Agilité de fréquence}
Un train non cohérent peut sauter de porteuse d'impulsion à impulsion sur une excursion $\Delta f$ ; la bande occupée
devient $B + \Delta f$ et le nombre d'échantillons indépendants de la cible vaut
$\min(n, \lfloor \Delta f/(c/2L)\rfloor + 1)$ pour une cible de longueur $L$, ce qui rend une cible déclarée Swerling~1/3
semblable à un cas 2/4 sur le train.

\subsection{Multifaisceau en émission}\label{sec:multibeam}
En S, un passage peut servir $B_{\mathrm{tx}}$ faisceaux. Le récepteur numérique forme déjà de nombreux faisceaux de
réception par position (sans perte ; une charge de calcul). Trois modes d'émission sont modélisés :
\begin{itemize}
  \item \textbf{Superposition en phase seule} ($B_{\mathrm{tx}} = 2, 4, 8$) : chaque TRM émet la somme à enveloppe
        constante des faisceaux (PA saturés, PAE inchangée), puissance $P/B_{\mathrm{tx}}$ par faisceau, gain
        d'émission complet, perte de \SIlist{0.4;0.6;0.8}{dB} (puissance perdue en intermodulation de la somme à
        enveloppe constante).
  \item \textbf{Partition de l'ouverture} ($B_{\mathrm{tx}} = 2, 4$) : moitiés ou quadrants (section~\ref{sec:subap}),
        puissance $P/B_{\mathrm{tx}}$ et gain $G\,N_{\mathrm{sub}}/N$ par faisceau, un faisceau plus large sur sa propre
        grille : moins de positions.
  \item \textbf{Multiplexage en temps} ($B_{\mathrm{tx}} = 2, 4, 8$) : impulsions à pleine puissance émises à la suite
        dans la PRI ; distance aveugle, réserve, facteur de charge et PAE calculés sur $B_{\mathrm{tx}}\tau$. Non
        proposé à l'optimiseur : avec des positions entrelacées, il dépense par position la même énergie et le même
        temps qu'un seul faisceau.
\end{itemize}

\subsection{Espace de recherche de veille}
L'espace de recherche S est le produit de : intégration (cohérente, non cohérente), mode d'émission et nombre de
faisceaux, largeurs de signal \SIlist{10;20;30;50;75;100;150;200}{MHz} compatibles avec la résolution du secteur,
impulsions de \SI{5}{\micro s} à \SI{2}{ms} (13 valeurs, la dernière fixée par la réserve), 1 à 32 impulsions, la
variante agile des trains non cohérents et, là où une surface encombre le secteur, des pondérations Doppler de Taylor
des trains cohérents. La PRI est la plus courte non ambiguë (ou \SI{1}{ms} pour les impulsions seules et les trains
ambigus autorisés). Environ \num{6500} candidats par secteur.

%=====================================================================================================================
\section{Traitement du signal}\label{sec:proc}
%=====================================================================================================================

\subsection{Blocs cohérents et cohérence}
Un train cohérent de $N$ impulsions est découpé en $M$ blocs ($M$ divise $N$) intégrés de façon non cohérente, $M$
minimisant le SNR requis de $M$ blocs plus la perte de cohérence d'un bloc. Les trains non cohérents, et les cas
Swerling~2/4 déclarés, ont $M = N$. La perte de cohérence d'un bloc de durée $T$ combine
\begin{itemize}
  \item la décorrélation de la cible $\eta = \exp(-(T/T_0)^2)$, avec $T_0 = \lambda/(2L\omega)$ tiré de la classe de
        cible (longueur $L$, vitesse de rotation $\omega$) ;
  \item la variance de phase résiduelle $\sigma^2$, qui donne $4,34\,\sigma^2\dB$ : bruit de phase de l'oscillateur
        filtré par l'aller-retour, $4\sin^2(\pi f\tau_{rt})$, intégré de $1/T$ à $B/2$ (un oscillateur pour l'émission
        et la réception ; un oscillateur photonique optimiste, \SI{-150}{dBc/Hz} à \SI{1}{kHz} de \SI{10}{GHz}), plus le
        mouvement résiduel de la plateforme $(4\pi\delta/\lambda)^2$ avec $\delta = \SI{10}{\micro m}$ RMS.
\end{itemize}
Les pertes sont négligeables avec ces valeurs ; l'effet visible est le plancher Doppler, donné par visée.

\subsection{Bancs de vitesses et d'accélérations}
La migration en distance est compensée (type keystone ou Radon, résidu \SI{0.1}{dB} quand la cible traverse plus de
1,5 case). Le filtre pour un écart de vitesse $\Delta v$ perd le terme intra-impulsion plus la migration sur la visée ;
le pas du banc est choisi pour qu'une cible à un demi-pas perde \SI{1}{dB} ; la perte \emph{moyenne} sur le pas entre
dans le bilan, et le nombre d'hypothèses dans le décompte des fausses alarmes et la charge de calcul. Les accélérations
jusqu'au maximum de la classe ajoutent une phase quadratique sur un bloc cohérent et une courbure résiduelle en
distance ; quand elles coûtent plus de \SI{1}{dB}, un banc d'hypothèses d'accélération multiplie les hypothèses de
vitesse.

\subsection{Grilles, CFAR et fausses alarmes}
La grille de détection est suréchantillonnée $\times4$ en distance et la transformée Doppler complétée de zéros
$\times4$ : chevauchement moyen \SI{0.05}{dB} et \SI{0.07}{dB}. Le CA-CFAR utilise \num{1024} cases de référence (perte
\SI{0.04}{dB}, exacte pour Swerling~1 sur un échantillon) ; l'OS-CFAR (rang 0,75) et le GO-CFAR (la plus grande des deux
demi-fenêtres, qui tient la $\Pfa$ à un bord de fouillis, \SI{0.05}{dB}) sont des options. Cases indépendantes par
position : étendue en distance sur $c/2B$, cases Doppler d'un bloc, hypothèses de vitesse et d'accélération, et cases
angulaires (faisceaux de réception par position $\times$ le carré de l'écart entre faisceaux de réception). Les grilles
suréchantillonnées ne sont pas comptées (échantillons corrélés).

%=====================================================================================================================
\section{Critère de détection}\label{sec:det}
%=====================================================================================================================

\subsection{Lois de détection exactes}
Un détecteur quadratique sur $n$ échantillons intégrés de façon non cohérente (les blocs cohérents) est modélisé
exactement : sous le bruit seul, la statistique suit une loi Gamma$(n,1)$, donc $\Pfa = P(\mathrm{Poisson}(T) \le n-1)$ ;
pour une cible stable de SNR total $\lambda$, $\Pd = \sum_k \mathrm{Pois}(k;\lambda)\,P(\mathrm{Poisson}(T)\le n+k-1)$.
Les cibles fluctuantes sont des mélanges sur une puissance de cible de loi Gamma de forme $k$ ($k=1$ : Swerling~1,
$2$ : Swerling~3, $n$ : Swerling~2, $2n$ : Swerling~4, $\infty$ : stable) ; les formes non entières sont admises, ce qui
fait entrer la décorrélation dans une visée et la diversité de polarisation. Les lois valent pour toute $\Pfa$ et
concordent avec l'ajustement de Shnidman à son erreur de \SI{0.3}{dB} près.

\subsection{Pd moyennée sur la position de la cible}
La position aléatoire de la cible coûte du SNR. Les écarts dans la maille hexagonale d'émission (diagramme réel), dans
la grille des faisceaux de réception, en chevauchement distance et Doppler, et au filtre de vitesse le plus proche
forment une loi de perte discrète (16 quantiles chacun, combinés comme indépendants). La $\Pd$ d'une visée est la $\Pd$
\emph{moyenne} sur cette loi, et le SNR qu'elle coûte est la perte de bilan appelée « positions ». Un pas de 1,4
largeur de faisceau d'émission est l'optimum de ce critère pour une visée (plat à \SI{0.1}{dB} près de 1,3 à 1,6) ; le
critère cumulé baissant la $\Pd$ demandée à une visée, l'optimiseur choisit le pas de chaque secteur de 1,3 à 2,0 par
pas de 0,1 (section~\ref{sec:opt}). Un coin de maille est à \SI{-7.9}{dB} au pas 1,4 : l'ordonnanceur doit décaler la
grille à chaque trame.

\subsection{Pd cumulée de la veille}
Une cible qui se rapproche à $v$ est vue une fois par revisite $RR$, aux distances $R(1+(k+\phi)q)$ avec $q = vRR/R$ et
une phase uniforme $\phi$ ; le SNR d'une visée varie en $(R/R_k)^4$. L'exigence porte sur la probabilité
\emph{cumulée}
\[
  P_{\mathrm{cum}}(R) = \bigl\langle 1 - \textstyle\prod_k (1 - P_{\mathrm{d},k})\bigr\rangle_\phi = 0,9
\]
à la portée exigée du secteur (les visées sous $\Pd = 10^{-3}$ ne comptent plus). Par rapport à une seule visée à $R$,
le SNR requis change de $+6,0\dB$ pour $q=1$ et de $-6,8\dB$ pour $q=0,1$.

\subsection{Politique de fausses alarmes}
Par défaut, un secteur tolère un nombre de fausses alarmes par revisite (\code{fa\_per\_rr} = 10) ; la $\Pfa$ par case
est ce nombre divisé par les positions et les cases indépendantes, bornée à $[10^{-14}, 10^{-4}]$ et arrondie au
vingtième de décade. Chaque fausse alarme est ensuite confirmée en Ku comme une détection. Les visées Ku utilisent
$\Pfa = 10^{-8}$.

%=====================================================================================================================
\section{Modèles de cibles}\label{sec:targets}
%=====================================================================================================================
Les cibles sont décrites par des \emph{classes} dont les valeurs sont des \emph{moyennes de classe} (l'équipement de
l'adversaire est inconnu). Un secteur nomme une classe et peut surcharger chaque valeur. ST est la cible de référence
conventionnelle du banc ; GENERIC complète un secteur qui ne nomme aucune classe.

\begin{center}
{\small\setlength{\tabcolsep}{4pt}\begin{tabular}{@{}lcccccccc@{}}
\toprule
Classe & $v$ (km/s) & $L$ (m) & $\omega$ (\si{\degree/s}) & $a_{\max}$ & SER S / Ku (\si{m^2}) & Fluct. & Dépol. & $T_0$ S / Ku (s)\\
\midrule
ST (référence) & 5 & 10 & 0,1 & 0 & \num{1000} / \num{1000} & SW1/2 & \SI{-10}{dB} & 3,1 / 0,6\\
GENERIC & 10 & 3 & 1 & 1\,g & 1 / 1 & SW1/2 & \SI{-10}{dB} & 1,0 / 0,20\\
KKV & 25 & 1,5 & 1 & 10\,g & 1,27 / 13,8 & SW3/4 & \SI{-20}{dB} & 2,0 / 0,41\\
SHIP & 5 & 100 & 0,1 & 1\,g & 300 / 300 & SW1/2 & \SI{-6}{dB} & 0,31 / 0,06\\
DEBRIS & 15 & 0,5 & 10 & 0 & 0,01 / 0,01 & SW1/2 & \SI{-3}{dB} & 0,61 / 0,12\\
DECOY & 25 & 1,5 & 0,5 & 0 & 0,135 / 0,135 & SW0 & \SI{-20}{dB} & 4,1 / 0,82\\
\bottomrule
\end{tabular}}
\end{center}

\paragraph{SER frontale et optique d'autodirecteur.} La SER est copolaire et frontale (l'aspect d'une cible qui se
rapproche). L'avant du KKV est un grand miroir d'autodirecteur infrarouge, métallique et donc conducteur parfait aux
fréquences radar, pointé vers sa cible. Sa contribution vaut
$\sigma_{\mathrm{opt}} = \max\bigl(\pi(2f)^2,\;\eta_r\,4\pi A^2/\lambda^2\bigr)$ avec $A = \pi D^2/4$,
$D = \SI{0.3}{m}$, $f/D = 1$, $\eta_r = 0,1$ : réflexion directe sur le primaire concave, et rétroréflexion « \oe{}il de
chat » de l'ensemble focal, en $1/\lambda^2$, donc surtout en Ku. Le leurre imite le corps sans l'optique, ce qui aide
la discrimination.

\paragraph{Fluctuation.} La fluctuation d'une visée est une forme de loi Gamma : la famille (SW1/2 : un degré de
liberté par échantillon indépendant, SW3/4 : deux, SW0 : stable) fois $\min(n, 1 + T_{\mathrm{obs}}/T_0)$ échantillons
indépendants sur la visée ($n$ blocs cohérents), fois la diversité de polarisation $(1+\rho)^2/(1+\rho^2)$ des deux
voies de réception. Un cas de Swerling déclaré prime sur la classe.

\paragraph{Masquage.} La scène est inconnue : le masquage d'une cible faible par une voisine forte est une sortie,
l'enveloppe de la réponse de la cible forte à 1, 2 et 3 cases en distance (avec la pondération de la visée) et en
Doppler (\SI{-13.3}{dB} au-delà d'une case sans fenêtre Doppler).

%=====================================================================================================================
\section{Environnement}\label{sec:env}
%=====================================================================================================================

\subsection{Structure du vaisseau}
Les obstacles sont dans le champ proche de l'ouverture (limite de Fraunhofer environ \SI{325}{m} en S) : le faisceau
est un tube de la largeur de l'ouverture. Des rayons partent de 37 points de l'ouverture dans la direction de pointage,
exactement sur les triangles du modèle 3D ; la fraction occultée $f$ coûte $-20\log_{10}(1-f)$ par passage, et une
position au-delà de \SI{3}{dB} par passage est masquée. La structure rayonne aussi : son émissivité fois sa
température physique (coque $\varepsilon = 0,01$ à \SI{290}{K}, radiateurs $0,05$ à \SI{750}{K}) est captée par le lobe
principal à travers la fraction occultée et par les lobes lointains comme un corps du ciel. Avec l'inclinaison de
\SI{25}{\degree}, la structure n'occulte que les directions à plus de \SI{71.7}{\degree} d'une normale de face, toutes
servies par une autre face à un dépointage plus faible : aucune position n'est masquée et les radiateurs ne sont jamais
vus ; la coque ajoute \SI{2.9}{K}, négligeable.

\subsection{Surfaces : ombre, fouillis, saturation du convertisseur}
Les corps à surface (Terre, Lune) sont des sphères à une distance $D$ donnée par leur rayon apparent. Une direction
dans le disque rencontre la surface à $d = D\cos\delta - \sqrt{R^2 - D^2\sin^2\delta}$ ; rien n'existe au-delà. La
rétrodiffusion de la surface est à « $\gamma$ constant », $\sigma_0 = \gamma\cos\theta_i$ ($\gamma$ = albédo radar ;
Terre \SI{-12}{dB} en S, Lune \SI{-11.5}{dB}). Le fouillis par les lobes secondaires dans la case distance $r$ de la
cible vient de l'anneau isodistance de la surface (aire $2\pi R r\,\Delta r/D$ par case), échantillonné sur 720 points
et vu à travers le diagramme aller-retour : $y = \sigma_c/\sigma_{\mathrm{cible}}$. L'écho d'une surface dans le lobe
principal est contrôlé contre le niveau d'écrêtage du CAN par élément (alerte), et un limbe dans un secteur est signalé
(GO-CFAR conseillé).

\subsection{Fouillis par forme d'onde}\label{sec:clutter}
En S, le fouillis est évalué \emph{par forme d'onde candidate et par tranche de dépointage de \SI{1}{\degree}}, hors du
budget de $\Tsys$ (la Ku garde un modèle plus simple de facteur de bruit) :
\begin{itemize}
  \item \textbf{Rapport de fouillis} $y$, par pointage et par faisceau d'émission de l'espace de recherche (face
        complète, ou diagramme 2D d'une sous-ouverture de partition), à la distance exigée du pointage, pour une case
        distance de référence ; il varie exactement en $c/2B$ avec la bande du candidat et est divisé par le nombre de
        blocs cohérents $M$ (l'exigence par bloc).
  \item \textbf{Résidu Doppler} $\rho(v)$ : le spectre Doppler de la surface (gaussien, moyenne et écart-type de sa
        vitesse de rapprochement donnés par le scénario), replié sur la PRF, à travers le filtre d'un bloc cohérent
        (poids uniformes $w_k$, ou Taylor \SIlist{-30;-40}{dB}),
        \[ H(x) = \Bigl|\textstyle\sum_k w_k e^{j2\pi kx}\Bigr|^2 \big/ \bigl(\textstyle\sum_k w_k\bigr)^2, \]
        plus un plancher blanc dû à la variance de phase résiduelle à travers la bande de bruit du filtre. $\rho = 1$
        sans blocs cohérents ou sans vitesse de rapprochement du scénario (prudent). Il est évalué à 11 vitesses de
        cible sur la plage du secteur.
  \item \textbf{Portée} : le fouillis est traité comme du bruit au détecteur, $S/(N + C)$ avec $C/N = a\,S/N$,
        $a = y\,\rho(v)/M$. Dans chaque tranche, $y$ est pris au quantile 99\,\% des pointages ; l'augmentation de SNR
        $G$ qui rétablit la $\Pd$ de la visée \emph{moyennée sur les vitesses de cible} est résolue par dichotomie, et la
        portée devient $R_n G^{-1/4}$. Si aucune puissance ne la rétablit, le candidat n'est pas jouable dans la tranche ;
        un pointage qu'aucun candidat ne dégage est aveuglé.
  \item \textbf{Exigence ramenée à la surface} : là où un pointage rencontre une surface à $d$ en deçà de la portée
        exigée, l'exigence devient $d - \text{garde}$ (garde = max(\SI{10}{km}, résolution en distance)), arrondie
        \emph{vers le haut} sur une grille de 5\,\% de la portée ; le critère cumulé est recalculé pour chacune de ces
        classes ; avec $d - \text{garde} \le 0$ il n'y a pas de volume et le pointage n'est pas balayé.
\end{itemize}

\subsection{Brouilleurs}
\paragraph{Les brouilleurs à distance} sont déclarés dans le ciel. Leur bruit entre dans $\Tsys$ à travers le
diagramme de réception : $T_j = (\text{densité de PIRE}/4\pi R^2)/k \times A_{\mathrm{eff}}(\text{direction})$. Un
brouilleur à plus d'une largeur de faisceau de réception de l'axe est dans les lobes secondaires et annulé
adaptativement de \SI{40}{dB} (les 8 plus forts par face).

\paragraph{Les brouilleurs d'autoprotection} sont portés par la cible du secteur (barrage sur la bande de leur classe,
densités de PIRE de \SI{5}{dBW/MHz} sur un KKV à \SI{27}{dBW/MHz} sur un vaisseau ; dans le lobe principal, non
annulés). Ils entrent exactement dans la portée : avec $S \propto 1/R^4$ et $J \propto 1/R^2$,
$S = \SNR_{\mathrm{req}}(N + J)$ est quadratique en $1/R^2$,
\[
  \frac{1}{R^2} = \frac12\left(\frac{1}{R_{bt}^2} + \sqrt{\frac{1}{R_{bt}^4} + \frac{4}{R_n^4}}\right),\qquad
  R_{bt} = \sqrt{\frac{P\,G_t\,\sigma\,\tau N}{4\pi\,J_0\,\SNR_{\mathrm{req}}\,L}},
\]
avec $R_n$ la portée contre le bruit seul et $R_{bt}$ la portée de percée. \textbf{Concept d'emploi} : la veille n'a
pas à percer un brouilleur d'autoprotection ; elle identifie l'axe du brouillage, et l'ordonnanceur pointe ensuite une
veille de percée Ku dans cet axe (gain d'émission Ku environ 25 fois celui de la S, donc une portée de percée environ
5 fois plus grande pour la même énergie). Un secteur dont le brouilleur met l'exigence hors de portée de la veille est
donc planifié contre le bruit seul et signalé pour la percée Ku.

\subsection{Autres termes}
Les objets non menaçants (débris, satellites) vus par revisite sont confirmés comme des cibles. La rotation de Faraday
sur l'aller-retour, $\Omega = 2\times2,36\cdot10^4\,B_\parallel\,\mathrm{TEC}/f^2$ (unités SI), est donnée à partir du
contenu électronique et du champ du scénario ; elle ne coûte rien avec le récepteur bipolarisation.

%=====================================================================================================================
\section{Équation radar et bilan des pertes}\label{sec:solve}
%=====================================================================================================================
La portée d'une visée est la forme énergétique de l'équation radar,
\[
  R^4 = \frac{P_{\mathrm{beam}}\,G_{\mathrm{TX}}\,G_{\mathrm{RX}}\,\lambda^2\,\sigma\,\tau N}
             {(4\pi)^3\,k\,\Tsys\,L\,\SNR_{\mathrm{req}}},
\]
avec $\SNR_{\mathrm{req}}$ le SNR sur l'énergie totale intégrée tiré de la loi de détection exacte (cumulée en veille),
$P_{\mathrm{beam}}$ et $G_{\mathrm{TX}}$ ceux du mode d'émission, $G_{\mathrm{RX}}$ avec la pondération de réception du
volume, et $\Tsys$ la température de budget de la direction. La perte totale $L$ est la somme de :
\begin{center}
\begin{tabularx}{\linewidth}{@{}lL@{}}
\toprule
Fixes (S) & précision finie \SI{0.3}{dB}, erreurs du réseau après étalonnage \SI{0.006}{dB} (section~\ref{sec:cal}),
occultation par les éléments Ku \SI{0.2}{dB}, chaîne d'émission \SI{0.6}{dB}\\
Bande & dégradation de bord moyennée sur le spectre occupé\\
Forme d'onde & compression (pondération, désadaptation NLFM, Doppler LFM), éclipse des trains décalés, chute de la
réserve, superposition multifaisceau, pondération Doppler\\
Antenne & désadaptation active, émission et réception\\
Polarisation & recombinaison GLRT ; énergie en polarisation croisée de la cible (un gain)\\
Traitement & positions (cases de grille, chevauchement, écart au filtre de vitesse, par la $\Pd$ moyennée), CFAR,
cohérence, accélération, résidu de migration\\
\bottomrule
\end{tabularx}
\end{center}
Pour la forme d'onde de référence (24 impulsions de \SI{200}{\micro s}, superposition de 4 faisceaux) sur la cible de
référence à \SI{2.8}{GHz} dans l'axe, les pertes totalisent environ \SI{6.6}{dB}, terme de positions compris (la
section~\ref{sec:effects} donne chaque terme et sa plage) ; une dérivation manuelle indépendante de l'équation concorde
à \SI{0.5}{\percent} (autotest).

Le coût d'une position de veille est électrique : énergie $P_{\mathrm{pk}}\,N\tau\,t/c_d/\eta$ et temps d'émission
$N\tau\,t/c_d$, avec $t$ les trains par passage (les faisceaux d'un multiplexage en temps, 1 sinon) et $c_d$ les
positions que couvre un passage (1, $B_{\mathrm{tx}}$, ou $B_{\mathrm{tx}}$ fois le rapport de maille de la
sous-ouverture).

%=====================================================================================================================
\section{Optimiseur de veille}\label{sec:opt}
%=====================================================================================================================
Pour chaque secteur, l'optimiseur (\code{core/optimizer.py}) choisit la fréquence centrale S et les formes d'onde, et
calcule l'énergie et le temps que chaque face y consacre.

\subsection{Fréquence}
Les fréquences centrales sont \emph{dérivées de la bande} : le bord bas utilisable plus des multiples de
\SI{5}{MHz}, gardées là où le signal tient dans la bande et ne touche aucune exclusion de spectre de la mission. La
recherche se fait en deux passes : les centres valides au pas de \SI{20}{MHz} (plus le plus central de chaque fragment
libre de la bande), puis le pas de \SI{5}{MHz} à \SI{\pm20}{MHz} du meilleur. L'optimum est plat en fréquence : la
grille sert surtout à éviter le spectre brouillé ou interdit. Un secteur peut imposer ses fréquences candidates.

\subsection{Plan de formes d'onde}
À chaque fréquence, les portées de tous les candidats sont calculées dans chaque tranche de dépointage de
\SI{1}{\degree}. Dans le modèle, le dépointage n'entre dans une portée que par les gains et la désadaptation active,
donc $R(\theta) = R(0)\,h(f,\theta)$ avec $h$ indépendant de la forme d'onde (vérifié à l'exécution) : une solution par
candidat et par fréquence suffit.
\begin{itemize}
  \item Une forme d'onde par groupe de tranches de \SI{5}{\degree}, au plus quatre formes d'onde distinctes par
        secteur ; les groupes voisins en partagent une par programmation dynamique sur des segments contigus.
  \item Un candidat est admissible sur un segment s'il atteint l'exigence dans chaque tranche peuplée. « Pas
        d'énergie en trop » : là où il dépasse l'exigence, l'\emph{impulsion est raccourcie} à $s\tau$ avec
        $s = (R_{\mathrm{req}}/R)^4$ (PA toujours saturés ; PAE exacte de l'impulsion raccourcie ; plancher
        $\tau B \ge 1$ ; les pertes de l'impulsion complète sont gardées, ce qui est prudent). Énergie et temps
        d'émission baissent tous deux.
  \item Objectif : minimiser $\max(\text{part d'énergie}, \text{part de temps d'émission})$ sur les faces, par une
        somme pondérée dont le poids est balayé (21 valeurs).
  \item Sans candidat admissible sur un segment : le moins de positions en deçà de la portée, puis la plus longue
        portée, puis le coût. Un secteur dont l'exigence est mise hors de portée par un brouilleur d'autoprotection est
        planifié contre le bruit seul (section~\ref{sec:env}).
  \item Classement des fréquences : d'abord la part non balayée due à la fréquence, puis les positions en deçà de la
        portée (tolérance 1\,\%), puis l'objectif, puis un critère secondaire déterministe (énergie, temps, nombre de
        formes d'onde, distance au centre de la bande c\oe{}ur).
\end{itemize}
Les candidats dominés (même fonction de coût, portée inférieure ou égale partout) sont élagués exactement avant le
plan. Chaque pas de grille de veille (1,3 à 2,0 largeurs de faisceau d'émission) est évalué à chaque fréquence
centrale, et le plan le moins cher est gardé : le pas change les positions, les pertes de grille, les faisceaux de
réception par position, les cases de fausses alarmes et le budget Ku qui suit. Quand une face est sursouscrite, la
revisite est allongée en proportion.

\subsection{Sorties et alertes}
Par secteur : fréquence et bande occupée, formes d'onde par tranche avec leur chronologie, portées au meilleur et au
pire dépointage, positions, faisceaux de réception par position, résumés du ciel et de l'environnement, énergie et
temps par face, revisite seul et dans la mission. Alertes : résolution non tenue, portée non atteinte, puissance ou
temps sursouscrits, part non balayée, désadaptation active, option de lobes secondaires de réception non tenue,
aveuglement ou montée du ciel au-delà du budget, masquage, limitation par le fouillis, surfaces plus proches que
l'exigence, brouillage d'autoprotection, revisite inférieure à trois temps de cohérence de la cible (le critère cumulé
suppose des visées indépendantes), exclusions de spectre retirant la plupart des centres.

%=====================================================================================================================
\section{Confirmation et poursuite active en Ku}\label{sec:ku}
%=====================================================================================================================
Chaque détection S, vraie cible ou fausse alarme, est confirmée en Ku (\code{core/tracking.py}, étape 3.2). Trois
modes :
\begin{itemize}
  \item \textbf{KU1+7} : une visée Ku sur l'estimation S, puis une grappe de 7 faisceaux en cas d'échec (deux
        allers-retours) ;
  \item \textbf{KU7} : la grappe de 7 faisceaux d'emblée (un aller-retour) ;
  \item \textbf{S2} : un « toucher » S pointé (une forme d'onde S de rôle CONFIRM, qui rejette aussi les fausses
        alarmes S), puis KU1+7 sur l'estimation angulaire améliorée.
\end{itemize}
La grappe est émise en 7 visées successives, en superposition en phase seule des 7 faisceaux (puissance / 7, perte
interpolée dans la table multifaisceau, \SI{0.76}{dB}), ou par une partition de la face en 2 ou 4 sous-ouvertures Ku ;
la moins chère est gardée. Les échos qui franchissent le seuil ont un SNR aléatoire, tiré par quantiles de la loi Gamma
de la cible ; l'erreur angulaire par axe vaut $K\,\mathrm{HPBW}/\sqrt{\mathrm{SNR}}$ (borne de Cramér--Rao $\times 1,13$)
plus un biais résiduel d'alignement S/Ku, moyennée par des n\oe{}uds de Gauss--Hermite. Chaque mode reçoit la plus petite
marge de puissance qui confirme une vraie cible avec une probabilité d'au moins 0,98, résolue par dichotomie ; le mode
qui minimise la charge de la face la plus chargée est choisi, ou l'opérateur en impose un.

La fréquence Ku est choisie par secteur sur une grille de \SI{10}{MHz} dérivée de la bande Ku, le signal dans la bande
et hors des exclusions Ku de la mission ; la cinquième harmonique S sur le canal Ku de la même face est signalée comme
règle de temps de l'ordonnanceur. Les visées Ku sont choisies dans leur propre espace (impulsions de
\SI{5}{\micro s} à \SI{2.08}{ms}, PRI de \SIrange{0.5}{5}{ms}, bandes de \SIrange{25}{200}{MHz}, les deux intégrations,
une variante agile des trains non cohérents), raccourcies à PA saturés juste assez pour la portée, avec le budget du
ciel Ku du secteur, son brouilleur d'autoprotection et le fouillis de ses surfaces par les lobes secondaires ; une visée
de percée Ku est pointée dans l'axe du brouillage quand le brouilleur met le secteur hors de portée de la veille.

La poursuite active est décrite par tranche de distance (nombre de pistes, période de mise à jour voulue par
l'opérateur, zone angulaire, précisions en distance et en vitesse radiale) et suit un modèle de piste
(\code{core/trackmodel.py}, étape 3.3). Chaque axe (distance ; deux axes transverses) porte un filtre de Kalman en
régime établi à trois états avec le modèle de man\oe{}uvre de Singer
($\sigma_m^2 = a_{\max}^2/3\,(1 + 4p_{\max} - p_0)$ plus un plancher, constante de temps par classe), discrétisé
exactement et résolu par l'équation de Riccati discrète. La zone angulaire ($\sigma$ = zone / 4) s'applique à l'angle
prédit à la mise à jour suivante, les précisions en distance et en vitesse à la piste filtrée ; chaque exigence devient
la plus grande erreur de mesure qui la tient, les planchers de mesure (\SI{0.05}{mrad}, \SI{0.05}{m}, \SI{0.01}{m/s})
retirés quadratiquement, et la précision en distance paie le résidu de couplage distance--Doppler
$c\tau\sigma_{f_d}/2B$. La visée de mise à jour est pointée sur la prédiction : son erreur de pointage entre dans sa
$\Pd$ comme une loi de perte sur le diagramme Ku réel ; elle couvre par des bancs $\pm 3\sigma$ de la vitesse radiale et
de l'accélération prédites ; la moins chère d'une visée pointée et d'une grappe séquentielle de 7 faisceaux est gardée.
Une chaîne de Markov sur les manques consécutifs (piste perdue après 3, l'erreur de pointage croissant pendant
l'extrapolation) donne le taux de perte ; les visées d'initialisation amènent une cible confirmée aux précisions
exigées ; les coûts sont moyennés sur 8 distances de la tranche. Quand la période de l'opérateur ne tient pas à une
distance, la plus longue période tenable est donnée en alerte. Un KKV est distingué d'un leurre sur les visées de
poursuite par quatre indices (dépolarisation, fluctuation, niveau de SER avec l'optique d'autodirecteur, man\oe{}uvre vue par
le filtre), l'incertitude de Faraday affaiblissant l'indice de dépolarisation sauf si la polarisation d'émission est
commutée. Un atténuateur après le LNA (\SI{30}{dB}) abaisse d'un facteur 5,6 la distance sous laquelle l'écho écrête les
CAN. S et Ku partagent l'énergie électrique d'une face ; leurs temps d'émission sont comptés séparément.

\paragraph{Réglages de calcul.} Depuis le 25 octobre 2026, les approximations sont admises quand elles ne dénaturent
pas le modèle, chacune une entrée \code{[tool.*]} dont la valeur exacte reste disponible : lois de perte de position
arrondies à \SI{0.02}{dB} (SNR requis à \SI{0.003}{dB} près), inversion cumulée à \SI{1e-3}{dB}, pas de grille cherché
à la fréquence retenue, trois centres S et trois centres Ku sur la bande c\oe{}ur quand le spectre est calme (aucune
exclusion, aucun brouilleur), marges S2 à \SI{0.05}{dB} avec 16 quantiles de SNR et 8 n\oe{}uds de Gauss--Hermite. Sur deux
missions de test, elles changent l'énergie de veille d'au plus \SI{+0.25}{\percent} (un secteur \SI{+9}{\percent}, un
\SI{-5}{\percent} : la fréquence compte pour certains secteurs) pour un calcul 7 à 8 fois plus rapide.

%=====================================================================================================================
\section{Ordonnanceur}\label{sec:sched}
%=====================================================================================================================
Le bilan des sections précédentes compte les ressources de chaque face \emph{en moyenne}. L'ordonnanceur
(\code{core/scheduler.py}, étape 4.3) vérifie que les tâches tiennent aussi \emph{dans le temps}, et combien de temps
une détection attend sa confirmation. S et Ku ont des chronologies d'émission distinctes sur chaque face.

\paragraph{Priorités.} Fixées par le concepteur :
\begin{center}
\begin{tabularx}{\linewidth}{@{}lL@{}}
\toprule
P0 & étalonnage (sans lui le radar ne fonctionne pas)\\
P1 & veille des volumes cochés \emph{autoprotection} (une case à cocher par volume)\\
P2 & confirmation d'une détection et initialisation de piste (ne jamais perdre une détection)\\
P3 & mises à jour de piste (quelques manques restent gérables) ; P3$'$ percée Ku, seulement si l'opérateur l'autorise\\
P4$_n$ & veille, avec une sous-priorité par volume choisie par l'opérateur ($n = 1$ servi d'abord)\\
P5 & réemploi de la marge : la part de veille ajoutée par une revisite raccourcie sous la valeur demandée
(\code{PRIO\_SEARCH}), ou de poursuite ajoutée par une cadence accélérée (\code{PRIO\_TRACK}) (section~\ref{sec:load})\\
\bottomrule
\end{tabularx}
\end{center}

\paragraph{Bilan statistique.} Sur chaque face et chaque bande, les tâches forment une file servie sans préemption
(une observation commencée n'est pas interrompue). L'attente moyenne de la classe $k$ suit la formule de Cobham,
\[
  W_k = \frac{W_0}{(1-\sigma_{k-1})(1-\sigma_k)},\qquad W_0 = \tfrac12\textstyle\sum_i \rho_i S_i,
\]
avec $\rho_i$ la part de temps de la classe $i$, $S_i$ sa durée de service et $\sigma_k$ la charge des classes jusqu'à
$k$ ; le quantile à 95\,\% utilise une queue exponentielle au-delà de la probabilité que l'émetteur soit occupé. Quand
la charge dépasse le temps, la veille est délestée d'abord (sa revisite s'étire, sous-priorité la plus basse d'abord),
puis les mises à jour de piste (leur période s'étire) ; l'étalonnage, l'autoprotection, la confirmation et
l'initialisation ne sont jamais délestés : s'ils ne tiennent pas, la mission n'est pas faisable.

\paragraph{Chronologie synthétique.} Une chronologie déterministe par face et par bande, sur la plus longue revisite
de la face (au plus \SI{60}{s}) : positions de veille réparties sur leur revisite, mises à jour de piste à leur
période, confirmations et initialisations tirées comme des arrivées de Poisson avec une graine fixe, un étalonnage ; à
chaque instant libre, la tâche prête de plus haute priorité, puis d'échéance la plus proche, est servie. Après
l'émission contiguë la plus longue que permettent le stockage de la face (environ \SI{0.3}{s}) ou la plaque froide, une
pause de recharge est insérée. La chronologie donne les retards de départ, la part des mises à jour en retard de plus
de \SI{10}{\percent} de leur période (comptées comme des manques), le retard de l'étalonnage et les pauses.

\paragraph{Ordre des pointages dans un volume.} Les pointages de veille d'un volume sur une face sont les points de sa
grille hexagonale en $(u,v)$ (où le faisceau garde une largeur constante, son élargissement aux grands dépointages est
donc exact), décalée d'une trame à la suivante de $0$, $(a+b)/3$ et $2(a+b)/3$ de la base de la grille (le décalage
d'une trame à l'autre que suppose le critère cumulé). Ils sont servis dans l'\emph{ordre de Morton à bits renversés}
(une suite de van der Corput à deux dimensions) : les $4^k$ premiers pointages tombent un dans chacun des $4^k$
sous-carrés du volume, à toutes les échelles ; la couverture est uniforme dès le début de la trame, et une trame étirée
ou coupée par le délestage se dégrade uniformément. Les cases du code ne sont pas plus fines que l'écart entre
pointages. Changer de pointage ne coûte rien avec le retard vrai numérique : aucune proximité n'est recherchée. Chaque
pointage dure le temps d'émission de la forme d'onde de sa tranche ; un passage multifaisceau sert ensemble les
pointages suivants. La chronologie synthétique sert ces passages ; l'outil donne par volume et par face les pointages,
le temps d'émission d'une trame, l'uniformité de la couverture à 10, 25 et 50\,\% de la trame, le pas moyen et la
cadence des pointages (entrées de l'étape de calcul) et la dispersion de la revisite des directions.

\paragraph{Règles entre faces et bandes.} Toutes les faces émettent une bande dans des créneaux communs (sinon
l'émission d'une face aveuglerait les récepteurs des autres) ; les créneaux sont dimensionnés par la face la plus
chargée, et toutes les faces y sont sourdes : les échos de veille perdent $-10\log_{10}(1-\rho_{\mathrm{créneaux}})$ ;
au-delà de \SI{0.5}{dB}, la perte est réinjectée dans le plan de veille par l'optimiseur global
(section~\ref{sec:load}). Sur une face où la cinquième harmonique d'une forme d'onde S couvre le canal Ku, les visées Ku
évitent les créneaux S.

\paragraph{Latence.} La latence d'une confirmation est l'attente à 95\,\% de la classe P2, plus son service, ses
allers-retours et la latence de traitement de la détection (section~\ref{sec:compute}). Elle est comparée à l'exigence
du secteur, ou par défaut à \SI{10}{\percent} du temps avant impact à la portée de détection.

%=====================================================================================================================
\section{Calcul et données}\label{sec:compute}
%=====================================================================================================================
Chaque face est entièrement numérique : chaque TRM a deux convertisseurs (un par polarisation) échantillonnant à
\SI{250}{Méch/s} sur 13 bits, soit \SI{63.7}{Tbit/s} par face en S et \SI{1.59}{Pbit/s} en Ku avant toute réduction.
Le bilan de calcul et de données (\code{core/compute.py}, étape 4.4) vérifie que la chaîne numérique suit ce que demande
le plan : débits, faisceaux de réception simultanés, charge de traitement, tampons, latence et puissance du traitement.
Chaque valeur ci-dessous est un réglage du fichier de paramètres (section \code{[compute]}).

\paragraph{Architecture.} Hiérarchique. Chaque TRSA ramène ses voies à la bande du signal (conversion numérique vers
le bas : échantillons complexes à $1,25\,B$, I et Q sur 16 bits). En S, une TRSA n'a que 50 voies, moins que les
faisceaux que forme un passage : elle envoie ses voies et le processeur de face forme les faisceaux. En Ku, une TRSA a
\num{1250} voies pour quelques faisceaux : elle forme des faisceaux partiels et les envoie (données divisées par environ
100). Le processeur de face forme ou somme les faisceaux, comprime les impulsions, exécute les bancs de vitesses et
d'accélérations et le CFAR, et annule les brouilleurs des lobes secondaires sur les 196 sorties de TRSA ; le calculateur
central tient les pistes.

\paragraph{Charge d'une visée.} Avec $f_s = 1,25\,B$, le temps d'écoute $T = N\cdot\mathrm{PRI}$,
$n_p = f_s\,\mathrm{PRI}$ échantillons par impulsion, $c_r = 4\,B\,\mathrm{PRI}$ cases distance de la grille de
détection, $h$ hypothèses de vitesse $\times$ accélération, $M$ cases Doppler et $b$ faisceaux de réception, les
opérations d'une visée sont
\begin{center}
\begin{tabularx}{\linewidth}{@{}lL@{}}
\toprule
formation de faisceaux & $8\times$ voies $\times\, b\, f_s T$ (une multiplication-accumulation complexe = 8 opérations)\\
compression d'impulsion & $10\, n_p\log_2 n_p$ par impulsion, faisceau et polarisation (convolution rapide)\\
bancs & $8\times 2\, b\, h\, c_r N$, plus la transformée Doppler $10\,M\log_2 M$ par case et hypothèse\\
CFAR & 10 par case testée et hypothèse\\
annulation des brouilleurs & $d^2$ par instantané d'apprentissage ($2d$ instantanés), $d^3/3$ par inversion, $d$ par
échantillon et faisceau\\
poursuite & $10^4$ par mise à jour de piste\\
\bottomrule
\end{tabularx}
\end{center}
Les visées sont les passages de veille S (cadence = positions / positions par passage / revisite), les confirmations Ku
et les mises à jour de piste Ku, tirées du plan et du budget Ku. Les fenêtres d'écoute d'une face se recouvrent ; par la
loi de Little, leur nombre moyen vaut $\sum \text{cadence}\times T$, et les faisceaux simultanés sont ce nombre arrondi
au-dessus fois les faisceaux du plus grand passage (une borne supérieure). Sur deux missions de test, des formes d'onde
S qui partitionnent l'ouverture en quatre sous-faisceaux d'émission sur 13,7 positions demandent 155 à 171 faisceaux de
réception simultanés ; le formateur forme 512 faisceaux en S et 32 en Ku (coût de calcul négligeable). Les tampons
cohérents sont contrôlés contre \SI{4}{To} de mémoire par face et le trafic mémoire (écritures des échantillons reçus
et comprimés, une lecture par passage : les bancs travaillent dans la mémoire interne du processeur) contre
\SI{50}{To/s}.

\paragraph{Puissance et capacité.} La puissance du traitement a un plafond (\SI{50}{kW} par face). Une part fixe
(\SI{20}{\percent} : fuites, horloges) est toujours consommée ; le reste suit la charge à \SI{0.01}{pJ} par opération et
\SI{0.05}{pJ} par bit transporté (2050--2060, bout optimiste). La capacité est ce que permet le plafond, environ
$4\times10^{18}$ opérations par seconde et par face ; les missions de test demandent $10^{11}$ à $2\times10^{16}$ (le
plus : les bancs de vitesses d'un secteur encombré par une surface), et le traitement consomme environ \SI{10}{kW} par
face. Les bilans de puissance et de chaleur qui dimensionnent le plan gardent le plafond (prudent : la consommation
dépend de la charge Ku, calculée après eux) ; la consommation réelle, l'appel du bus et la chaleur de face corrigés sont
donnés à côté, et la répartition de la puissance entre faces utilise la consommation réelle.

\paragraph{Latence.} Les données arrivent au processeur de face pendant la réception des échos ; après la fin d'un
passage restent le retard d'un lien plus lent que le flux (aucun aujourd'hui : \SI{0.2}{Tbit/s} par lien de TRSA contre
\SI{4}{Tbit/s}, \SI{39}{Tbit/s} vers le processeur de face contre \SI{400}{Tbit/s}) et le traitement à la capacité que
laisse la charge moyenne. Cette latence de traitement (\SI{0.01}{\micro s} à \SI{5}{ms} sur les missions de test) entre
dans la latence de confirmation de l'ordonnanceur.

\paragraph{En vol.} Le calculateur radar applique le plan de mission, l'ordonnanceur et la poursuite, et resélectionne
un plan de veille dans la \emph{table pour le vol} du secteur quand le spectre mesuré change. La table garde un plan
par fréquence centrale S évaluée par la recherche (son meilleur pas de grille), le plan retenu d'abord, au plus 8
répartis sur la bande ; un plan de couverture moindre ou avec plus de positions en deçà de la portée est gardé et
signalé dégradé. En vol, le calculateur prend la meilleure entrée dont la bande occupée est libre et dont la charge
tient sur la face (non dégradée d'abord), sinon il alerte et déleste par priorité. Une entrée coûte \SI{64}{o} par
groupe de tranches (quelques ko par secteur, dans la mémoire de la face à côté des tampons cohérents) et \num{1000}
opérations à réévaluer (provisoire). La réoptimisation complète reste en préparation de mission ; la fréquence Ku de
confirmation n'a pas encore de table pour le vol.

%=====================================================================================================================
\section{Re-pointage en vol}\label{sec:repoint}
%=====================================================================================================================
Les repères de la section~\ref{sec:frames} ramènent chaque volume dans le repère du vaisseau pour un message de
navigation. En vol, le vaisseau tourne (jusqu'à \SI{30}{\degree/s} et \SI{15}{\degree/s^2}, réglages éditables) et se
déplace sur son orbite ; le calculateur radar garde les pointages sur leurs volumes et leurs pistes
(\code{core/repoint.py}, étape 4.5).

\paragraph{Boucle.} La grille de pointage d'un volume est fixée une fois par trame dans le repère propre du volume ;
chaque pointage est converti en $(u,v)$ d'une face à l'instant de son émission, avec le dernier message de navigation
prolongé par la vitesse mesurée, et émis par la face qui le voit alors ; les faisceaux de réception sont formés avec
l'attitude de la réception. La couverture d'un volume reste donc exacte pendant que le vaisseau tourne au cours d'une
trame. Les volumes liés au vaisseau ne bougent pas ; le décalage du limbe suit l'altitude.

\paragraph{Erreurs.} À une émission, la navigation a au plus le retard plus un intervalle de message
(\SI{15}{ms}) ; l'erreur quadratique moyenne de pointage vaut
\[
  \overline{r^2} = 2\left(\sigma_{\mathrm{att}}^2 + (\sigma_{\dot\theta}\,a)^2\right) + \left(\tfrac12\,\alpha_{\max}a^2\right)^2,
\]
avec la connaissance d'attitude et de vitesse $\sigma_{\mathrm{att}} = 1''$, $\sigma_{\dot\theta} = 10^{-4}\,\si{\degree/s}$
par axe et l'accélération non mesurée $\alpha_{\max}$, qui domine ($6''$). Sa perte moyenne aller-retour de forme de
faisceau $2\times12,04\,\overline{r^2}/\theta_3^2$\,dB rejoint les pertes de chaque bande (\num{3e-5}\,dB en S,
\num{8e-4}\,dB en Ku) ; l'erreur est contrôlée contre \SI{5}{\percent} d'une largeur de faisceau (elle vaut $0,001$ en
S, $0,006$ en Ku). Dans une visée de poursuite, l'erreur par axe s'ajoute quadratiquement à l'erreur de pointage du
filtre. La vitesse angulaire elle-même est mesurée et compensée : elle n'entre pas dans l'erreur.

\paragraph{Migration entre faces.} Un volume non lié au vaisseau traverse une frontière de face à la vitesse maximale
en une fraction de seconde à environ une seconde ; la face voisine doit porter sa charge pendant ce temps (alerte
sinon). L'enveloppe complète des attitudes est contrôlée plus tard (moteur de validation).

\paragraph{Calcul.} Environ 100 opérations par pointage et, par trame, le masque de chaque volume non lié au vaisseau :
environ $10^7$ opérations par seconde et par face, négligeable.

%=====================================================================================================================
\section{Pilotage de la charge}\label{sec:load}
%=====================================================================================================================
Cette section rassemble les règles qui décident quelle tâche reçoit quelle ressource, quand, et ce qui cède quand une
ressource manque (étapes 4.1 à 4.5 et 5.2). Elles dimensionnent et vérifient un plan en préparation de mission ; leur
exécution dynamique en vol relève de la simulation globale.

\subsection{Ressources}
Chaque face a ses propres ressources, sauf l'énergie et le rejet de chaleur, partagés par le vaisseau.
\begin{center}
{\small\begin{tabularx}{\linewidth}{@{}lllL@{}}
\toprule
Ressource & Portée & Capacité & Charge\\
\midrule
Énergie électrique & vaisseau ; face & \SI{8}{MW} ; alimentation \SI{8}{MW}, plafond opérateur & appel moyen du bus / capacité\\
Stockage de face & face & \SI{5}{MJ} (\SI{80}{\percent} utilisables), \SI{15}{MW} & déficit par visée et par pire rafale\\
Temps d'émission S & face & 1 & somme des parts de temps des tâches S\\
Temps d'émission Ku & face & 1 & idem en Ku (réseau distinct)\\
Chaleur & face ; vaisseau & \SI{15.9}{MW} ; \SI{6}{MW} & chaleur moyenne / capacité\\
Calcul & face & $4\times10^{18}$ op/s & opérations par seconde / capacité\\
Faisceaux de réception & face, bande & 512 S, 32 Ku & pic de faisceaux d'un passage / capacité\\
Mémoire, débit & face & \SI{4}{To}, \SI{50}{To/s} & tampons cohérents au pic ; trafic\\
\bottomrule
\end{tabularx}}
\end{center}
La \emph{ressource dimensionnante} d'une face est celle dont la charge est la plus haute ; pour l'énergie, le plus grand
de (besoin du vaisseau / \SI{8}{MW}) et (besoin de la face / son plafond).

\subsection{Enchaînement du calcul d'une mission}
\begin{enumerate}
  \item \textbf{Plan de veille par secteur} (section~\ref{sec:opt}) : fréquence, grille, forme d'onde par tranche de
        dépointage, impulsion raccourcie à l'exigence ; objectif $\min\max$(part d'énergie, part de temps S) sur les
        faces, l'énergie rapportée à \SI{2}{MW} par face. Une direction est servie par la face la plus proche.
  \item \textbf{Budgets} : totaux par face, confirmation Ku (mode choisi pour baisser la face la plus chargée) et
        poursuite, calcul et données, répartition de la puissance (ci-dessous), re-pointage, ordonnanceur.
  \item \textbf{Boucle globale} (au plus 3 passes) : les couplages sont réinjectés dans les plans (éclipse des créneaux
        communs au-delà de \SI{0.5}{dB} ; coût Ku des fausses alarmes quand elles dépendent de la forme d'onde ; formes
        d'onde demandant plus de faisceaux que le formateur retirées) ; chaque secteur peut passer à l'une de ses
        alternatives (autre fréquence ou grille, même couverture et même portée) quand elle baisse la face la plus
        chargée de la mission, puis la somme des charges ; seuls les secteurs concernés sont replanifiés ; la boucle
        s'arrête quand rien ne change.
  \item \textbf{Réemploi de la marge} selon le mode de réallocation de la mission (ci-dessous), puis \textbf{table de
        saturation} de la poursuite.
  \item \textbf{Verdict} par ressource et par face, sur les revisites demandées.
\end{enumerate}

\subsection{Répartition de la puissance entre les faces}
Le besoin de chaque face (appel du bus) est décomposé par niveau de priorité : P0 permanent (avec la consommation réelle
du traitement) et étalonnage, P1 veille d'autoprotection, P2 confirmation et initialisation, P3 mises à jour de piste
par classe de cible (KKV, SHIP, DECOY, DEBRIS, autres), P3$'$ percée, P4$_n$ veille, P5 réemploi de la marge.
\begin{itemize}
  \item \textbf{Puissance suffisante} : chaque face reçoit une part des \SI{8}{MW} proportionnelle à son besoin ;
        l'excédent d'une face plafonnée va aux autres.
  \item \textbf{Manque} : les niveaux sont servis dans l'ordre, toutes faces confondues (une veille P1 d'une face avant
        une piste P3 d'une autre) ; le dernier niveau servi l'est en partie : ses volumes de veille reçoivent une
        revisite tenable ($RR$ / part servie), ses pistes un nombre tenable. Un niveau P0 à P2 non servi en entier rend
        la mission non tenable.
  \item \textbf{Verdict} : la charge d'énergie d'une face est son besoin divisé par son allocation, la même pour toutes
        les faces quand la puissance suffit (besoin total sur \SI{8}{MW}), au-dessus de 1 pour les faces dont des
        niveaux ne sont pas servis. L'alimentation et le plafond opérateur d'une face restent un contrôle à part.
\end{itemize}

\subsection{Réemploi de la marge}
Deux réglages de mission décident de la marge : le \textbf{mode de réallocation} \code{realloc\_mode} (\code{NONE} par
défaut : la marge reste libre ; \code{PRIO\_SEARCH} : elle raccourcit les revisites de veille ; \code{PRIO\_TRACK} : elle
accélère seulement la cadence de poursuite, les volumes restant à leur revisite demandée) et le \textbf{plafond de
charge} \code{load\_ceiling} $X$ (\SI{80}{\percent} par défaut), toujours actif : au-delà commence une réserve de
sécurité jamais consommée (entrelacement imparfait), et une face qui l'entame est signalée. En \code{PRIO\_SEARCH}, la
revisite demandée pour un volume est un minimum ; les revisites sont raccourcies jusqu'à ce que la ressource
dimensionnante atteigne $X$ :
\begin{itemize}
  \item un facteur $g \ge 1$ pour tous les volumes de veille, autoprotection comprise, monté ensemble ; les volumes
        d'une face qui atteint $X$ s'arrêtent, les autres continuent ; un volume sur plusieurs faces est borné par la
        plus chargée ; un volume s'arrête aussi à sa \code{rr\_min\_s}, \SI{1}{s} par défaut (\SI{0.1}{s} pour une
        tranche de poursuite en \code{PRIO\_TRACK}) ;
  \item le plan est gardé (formes d'onde, grille, durées d'impulsion) ; la revisite devient $RR/g$ ; les fausses
        alarmes restent au même nombre par seconde ($\Pfa$ par case divisée par $g$), donc la charge de confirmation ne
        bouge pas ; le gain est donné en portée de détection cumulée à la revisite réelle ;
  \item la charge de veille ajoutée est le niveau P5, le premier rendu en cas de surcharge ; le verdict reste porté sur
        la revisite demandée ;
  \item une sonde à $g = 2$ donne la pente de chaque ressource ; après le remplissage, les budgets sont réassemblés aux
        revisites réelles et une ressource non linéaire est corrigée par la sécante de ses deux derniers points (au plus
        2 corrections, tolérance 0,5 point) ; si une ressource reste au-dessus du plafond (les faisceaux simultanés,
        comptés en entiers), les facteurs sont ramenés ensemble par bissection.
\end{itemize}
Exemple (mission de test, $X = \SI{80}{\percent}$, plafonds opérateur de \SI{2}{MW} par face) : des revisites de 20,
5 et \SI{120}{s} deviennent 0,09, 0,068 et \SI{1.6}{s}, bornées par l'énergie des faces 1 et 2, portées cumulées
$+56$ à $+89\,\%$ ; avec \code{rr\_min\_s} = \SI{5}{s} sur le premier volume, $g = 4$ et $+19\,\%$ (essai fait avant
le plancher par défaut de \SI{1}{s}).

\subsection{Saturation de la poursuite}
La spécification de l'opérateur ($n$ pistes à la période $x$ par tranche, volumes à leur revisite $y$) est le point de
dimensionnement. Au-delà de $n$ :
\begin{center}
{\small\begin{tabularx}{\linewidth}{@{}lL@{}}
\toprule
État & Comportement\\
\midrule
Saturation gérable & $n+a$ pistes à $x$, veille à $y$ : la place jusqu'au plafond (marge réemployée comprise) divisée
par le coût d'une piste, le plus petit sur toutes les ressources ; permise aussi en \code{NONE}\\
Saturation non gérable & \code{PRIO\_SEARCH} et \code{NONE} : veille tenue à $y$, pistes allongées dans l'ordre strict
des classes (DEBRIS d'abord, KKV en dernier) ; \code{PRIO\_TRACK} : pistes tenues à $x$, veille dégradée dans l'ordre
P4$_n$, puis pistes allongées ; P0 et P1 intouchables ; quand le temps Ku limite, d'abord percée retirée et
confirmations basculées en S2\\
Saturation critique & planchers atteints (perte de piste par mise à jour de \SI{1}{\percent} ou latence, revisite maximale
d'un volume) : abandon de tâches entières selon le mode (volumes P4$_n$ puis pistes en \code{PRIO\_TRACK}, l'inverse
sinon), puis P2\\
P2 inatteignable & alerte majeure, le radar continue en P0 et P1 seulement, plus aucune création de piste\\
\bottomrule
\end{tabularx}}
\end{center}
En préparation de mission, le nombre de pistes est une entrée : l'outil publie, par tranche, les \textbf{seuils} de
ces états (pistes en plus gérables, seule et en cumul, entrée en saturation critique, premier abandon). Le coût d'une
piste vient d'une sonde (une piste de plus ou \SI{10}{\percent} de $n$) ; la valeur d'une tranche seule est une
extrapolation linéaire, le cumul (même fraction de débordement pour toutes les tranches) est vérifié par un assemblage.
Une piste allongée coûte son coût à $x$ multiplié par $x/T$ (approximation optimiste : la prédiction plus large n'est pas
recalculée).

\subsection{Surcharge : ordre de ce qui cède}
\begin{center}
{\small\begin{tabularx}{\linewidth}{@{}lL@{}}
\toprule
Ressource en manque & Réponse\\
\midrule
Temps S (plan de veille) & revisites de la face allongées en proportion (alerte)\\
Énergie & niveaux servis dans l'ordre ; revisite ou nombre de pistes tenable du dernier niveau\\
Temps (ordonnanceur) & P4$_n$ à partir de la sous-priorité la plus basse (revisite étirée), puis P3 (période de mise à
jour étirée) ; P0, P1, P2 jamais délestés (mission non faisable) ; en vol, P5 est rendu en premier\\
Faisceaux & formes d'onde demandant plus de faisceaux que le formateur retirées du plan\\
Migration entre faces & la face voisine doit porter la charge du volume (alerte)\\
\bottomrule
\end{tabularx}}
\end{center}
Avec l'ordre de Morton à bits renversés des pointages (section~\ref{sec:sched}), une trame étirée ou coupée garde une
couverture uniforme.

\subsection{En vol}
Le calculateur radar re-pointe chaque pointage à son émission (section~\ref{sec:repoint}) et resélectionne des formes
d'onde précalculées quand le spectre mesuré change. Règles fixées pour la simulation globale : la répartition devient
dynamique (ressources allouées et reprises entre fonctions selon la charge réelle ; le réemploi de la marge s'applique,
P5 rendu en premier) ; sur une surcharge, avant tout délestage, la ressource qui bloque la face en limite est
réallouée : énergie depuis les faces qui ont de la marge, temps d'émission par resélection de formes d'onde (énergie
contre temps) ou par la face voisine dans les zones de recouvrement, chaleur, calcul, faisceaux et mémoire par le
recouvrement ; le délestage par priorités vient en dernier.

\subsection{Verdict de la mission}\label{sec:verdict}
Le verdict (\code{core/validation.py}, étape 5.3) juge le plan en préparation de mission seulement, sur l'instantané de
navigation de la mission (un seul message de navigation : l'enveloppe d'attitudes relève du simulateur de scénario) et à
la spécification demandée (le réemploi de la marge est rapporté à part). Il n'ajoute aucune physique : chaque contrainte
est lue dans les sorties des autres modules comme une charge, besoin sur limite, de l'un de trois genres :
\begin{center}
{\small\begin{tabularx}{\linewidth}{@{}lLl@{}}
\toprule
Genre & Contraintes & États\\
\midrule
Ressource & énergie de la face (besoin sur allocation), temps S et Ku, chaleur, calcul, mémoire et son débit, faisceaux S
et Ku, lien TRSA, entrée du processeur, puissance du calcul, stockage par observation ; énergie et rejet de chaleur du
vaisseau & OK $\le$ \code{load\_ceiling} $<$ ALERTE $\le 1 <$ NOK\\
Exigence & plan trouvé, portée, revisite dans la mission, mode de confirmation admissible, probabilité de confirmation
(charge $(1-P)/(1-p_\mathrm{min})$), latence, période et portée de chaque tranche de poursuite, discrimination, niveau
de puissance servi en entier, ordonnancement faisable, délestage & OK $\le 1 <$ NOK\\
Transitoire & migration entre faces quand le vaisseau tourne, erreur de pointage de navigation, accélération angulaire
au-delà du dimensionnement & OK $\le 1 <$ ALERTE\\
\bottomrule
\end{tabularx}}
\end{center}
Des positions courtes dans la tolérance du secteur et une discrimination atteinte sous la plus courte portée d'intérêt
sont des alertes ; le stockage et la chaleur d'une revisite tassée sont des règles de l'ordonnanceur (la frise
synthétique alerte s'il reste une pause ou une mise à jour en retard). La mission est NOK dès qu'une contrainte est NOK, la
spécification de l'opérateur n'étant pas tenue ; chaque contrainte porte son niveau de priorité, et les niveaux en défaut
sont listés. Le verdict donne aussi la ressource limitante (la plus forte charge de ressource) de la mission et de chaque
face, et l'état de saturation à la spécification (nominal, réserve entamée, sur-souscrit ou P2 inatteignable) avec la
place par tranche de poursuite de la section~\ref{sec:sched}.

\paragraph{Conseil de relâchement.} Chaque contrainte en alerte ou en défaut reçoit ses leviers chiffrés. Une ressource
linéaire d'une face (énergie, temps S ou Ku) au-dessus de sa cible (le plafond de charge ; la valeur à 1 en second
chiffre) de $\Delta L$ est soulagée par un élément de part $s>\Delta L$ dans sa charge, qui garde la fraction
$f=1-\Delta L/s$ de sa charge. La portée d'un volume est la portée au pire : un objet raté en début de cycle doit encore
être détecté à $R$, les observations qui comptent sont donc espacées de $v\,RR$ au-delà de $R$ (pour 20~km/s et une
revisite de 20~s, un volume à 500~km doit détecter entre 500 et 900~km), et un volume de veille suit
\[ \frac{L'}{L}=\Big(\frac{R'}{R}\Big)^4\,\frac{\mathrm{SNR}_\mathrm{req}(q')}{\mathrm{SNR}_\mathrm{req}(q)}\,
\frac{RR}{RR'},\qquad q=\frac{v_\mathrm{close}\,RR}{R}, \]
le critère cumulé de la section~\ref{sec:det} sur sa forme d'onde principale, résolu pour sa portée $R'$ ou sa
revisite $RR'$ (une revisite plus longue espace les observations et peut alourdir la charge d'une cible rapide : pas de
levier de revisite alors), ou par sa SER de dimensionnement $\sigma'=\sigma/f$ (une hypothèse de menace) ; une tranche de poursuite par son nombre de pistes
$n'=nf$, sa période $T'=T/f$ ou la précision qui fixe l'énergie d'une mise à jour, $a'=a/\sqrt f$ (temps d'émission
pris limité en énergie, d'où les mêmes lois). Les trois plus gros éléments sont donnés seuls, plus un facteur uniforme
sur les revisites de toute la face quand aucun ne suffit seul ; les alternatives de confirmation (exactes), le plafond
opérateur de la face et la percée sont donnés tels quels. Quand la puissance du vaisseau suffit, l'allocation est au
prorata des besoins et la charge d'énergie de chaque face est celle du vaisseau : le conseil est donné une fois, sur le
vaisseau. Les ressources en marches (faisceaux, calcul, débit mémoire, stockage) et la chaleur reçoivent un facteur
proportionnel, signalé approché, sans la part statique de la puissance du calcul ; une crête par observation (lien
TRSA, entrée du processeur, tampons) est fixée par les formes d'onde, pas par les revisites, et n'a pas de levier de
revisite. Une exigence en défaut reçoit sa valeur directe (défaut accepté ou portée des tranches d'angle courtes,
probabilité de confirmation ou latence atteintes, plus longue période de poursuite tenable, revisite ou pistes tenables
d'un niveau de puissance servi en partie). Chaque chiffre est une estimation : la vérification recalcule la mission avec
le changement et procède par bissection (six recalculs) pour le plus petit relâchement qui tient, et seul le chiffre
vérifié fait foi.

\paragraph{Redondances et recouvrements.} Information seulement : le plan ne change pas. À la fin du calcul, le test
« le plan de A tient le besoin de B » utilise les plans retenus, avec la règle de la portée
au pire (revisite de A, cible et portée de B, critère cumulé) ; chaque volume redondant est donné avec son coût par
face, et une ressource de face en alerte ou en défaut (énergie, temps S) reçoit d'abord le conseil de le retirer (besoin
entièrement tenu) ou de le découper en boîtes restantes (besoin tenu en partie), spécification gardée ; la vérification
applique le retrait ou les morceaux. Une direction est servie par la face de moindre dépointage ; la face voisine la
voit aussi tant que son propre dépointage reste sous la limite. Par face et voisine, le verdict donne l'angle solide de
cette zone, le temps S et l'énergie de la face dans la zone et leur coût sur la voisine, au dépointage plus fort : coût
par position du plan là où il a cette tranche, au-delà forme d'onde de la tranche la plus proche avec sa portée mise à
l'échelle entre les deux angles, échelle d'impulsion $(R/r)^4$ et loi de coût de la veille (inatteignable au-delà de
l'impulsion pleine), avec la densité de positions de la voisine ; puis la marge de la voisine jusqu'au plafond de charge
(temps S ; énergie de la face si elle est plafonnée, sinon du vaisseau) et la part qu'elle pourrait reprendre. Les zones
de responsabilité restent celles du plan : la réallocation en vol est en phase 6. Le verdict porte sur les
valeurs de la menace données (celles du secteur, sinon celles de sa classe) ; il ne cherche pas de combien la menace peut
s'en écarter : pour une autre hypothèse, l'opérateur édite la classe ou surcharge le secteur et relance le calcul.

%=====================================================================================================================
\section{Réglages de mission et présentation des cibles}\label{sec:mission51}
%=====================================================================================================================
L'étape 5.1 fait des réglages opérationnels des clés de mission, chacune par défaut à la valeur du fichier de
paramètres : la probabilité de confirmer une vraie cible (par mission, ou par volume), la $P_d$ cumulée de veille et les
fausses alarmes tolérées (par revisite ou par seconde) d'un volume, la $P_d$ d'une mise à jour de piste et la
probabilité de discrimination d'une tranche de distance, et la plus grande vitesse et accélération angulaires du
vaisseau (la perte de bande de la navigation reste dimensionnée par le fichier ; une mission plus sévère est signalée
avec la perte qui manque). Un interrupteur de mission autorise la percée Ku dans l'axe d'un brouilleur
d'autoprotection : elle est alors budgétée (une visée par revisite) et ordonnancée après les mises à jour de piste
(P3$'$) ; elle est interdite par défaut. L'étape 5.2 ajoute \code{auto\_alloc\_max} par mission et \code{rr\_min\_s}
par volume (section~\ref{sec:load}).

\paragraph{SER et présentation.} Chaque classe de la bibliothèque a une table de sa SER en fonction de la présentation,
tous les \SI{30}{\degree} en azimut (12 valeurs) et en élévation (7 valeurs), en dB par rapport à la SER frontale,
interpolée sur une grille de \SI{2}{\degree} qui contient le nez. La détection utilise la \emph{pire} présentation
(moindre SER totale, l'optique d'autodirecteur comptée seulement dans son cône de \SI{\pm5}{\degree}), sauf pour le KKV,
pris \emph{de face} : il vise le vaisseau, nez et autodirecteur vers le radar. Avec les tables actuelles, aucune portée
de détection ne change. La dispersion de la table sur toutes les présentations (aucune pour une classe vue de face) et
la dispersion d'un objet à l'autre d'une classe (3 à \SI{5}{dB} en SER, \SI{2}{dB} en dépolarisation) entrent dans les
indices de discrimination comme des termes systématiques que des visées supplémentaires ne réduisent pas.

%=====================================================================================================================
\section{Autres fonctions}
%=====================================================================================================================
\paragraph{Erreurs du réseau et étalonnage.}\label{sec:cal} Chaque TRM est étalonné en vol par son coupleur (bouclage en
émission et en réception, tous les TRM d'une face à la fois avec des codes orthogonaux), \SI{10}{ms} en S et
\SI{50}{ms} en Ku toutes les \SI{300}{s}, pris sur le temps d'émission de la face. Les erreurs résiduelles en fin
d'intervalle (aléatoires par TRM \SI{1}{\degree} / \SI{0.1}{dB}, corrélées par TRSA \SI{0.5}{\degree} / \SI{0.05}{dB}
plus une dérive thermique de \SI{0.1}{\degree/min} / \SI{0.01}{dB/min}) donnent la perte de gain par la formule de Ruze
en émission et en réception (environ \SI{0.006}{dB}, à la place d'une perte globale de \SI{0.5}{dB} / \SI{0.8}{dB}), un
plancher moyen de lobes secondaires vers \SI{-59}{dB} et une erreur d'axe de quelques centièmes de milliradian ; le
plancher de mesure angulaire et le biais d'alignement S/Ku restent des entrées, contrôlées contre eux. Les modes
dégradés et les pannes de TRM sont traités en fin de développement (fonctionnement nominal d'abord).

\paragraph{Redondance.} Un secteur $A$ couvre le besoin d'un secteur $B$ si, avec les formes d'onde et la revisite de
$A$, la cible, la portée et la pondération de réception de $B$ (le récepteur numérique peut appliquer les deux sur la
même visée), le critère cumulé de $B$ est tenu, avec une PRI non ambiguë et une résolution suffisante.

\paragraph{Estimation rapide.} Pour un retour immédiat pendant que l'opérateur édite, une table précalculée donne
l'énergie la moins chère par position en fonction de la portée, de la SER, de la vitesse et du dépointage ; le verdict
vient du calcul complet.

\paragraph{Repères et navigation.}\label{sec:frames} Le message de navigation porte un quaternion de Hamilton (vaisseau
vers ICRF) et une orbite facultative. Avant tout calcul, la mission est ramenée une fois dans le repère du vaisseau :
chaque secteur donné dans un autre repère (inertiel, RTN, relatif au limbe, deux directions) reçoit sa matrice de
rotation et la trame des directions est tournée avant le test de forme ; le ciel inertiel est tourné en sens inverse.
Tous les calculs suivants restent dans le repère du vaisseau.

%=====================================================================================================================
\section{Logiciel}\label{sec:software}
%=====================================================================================================================

\subsection{Modules}
\begin{center}
\begin{tabularx}{\linewidth}{@{}>{\raggedright\arraybackslash}p{0.33\linewidth}L@{}}
\toprule
Module & Rôle\\
\midrule
\code{core/params.py} & charge \code{radar\_system.toml} une fois par processus ; un autre fichier par \code{RADAR\_PARAMS}\\
\code{core/radar\_model.py} & bandes, TRM, diagrammes, pertes, formes d'onde, traitement, cibles, \code{solve()}\\
\code{core/antenna.py} & diagrammes à partir de la composition, profils des lobes lointains, sous-ouvertures (en cache)\\
\code{core/detection.py} & lois de détection exactes, $\Pd$ moyennée et cumulée\\
\code{core/sky.py} & bruit du ciel par direction, aveuglement, budget du secteur\\
\code{core/environment.py} & structure (modèle 3D, lancer de rayons, en cache), surfaces et fouillis, résidu Doppler,
brouilleurs, Faraday\\
\code{core/isolation.py} & fuites entre bandes et entre faces, règles d'ordonnanceur\\
\code{core/frames.py} & navigation, repères des volumes\\
\code{core/optimizer.py} & optimiseur de veille S, calcul de mission, boucle globale, résumés des secteurs\\
\code{core/tracking.py} & budget de confirmation Ku et de poursuite active\\
\code{core/trackmodel.py} & modèle de piste : filtre de Kalman de Singer, chaîne de perte, initialisation, discrimination\\
\code{core/scheduler.py} & priorités, attentes, chronologie synthétique, créneaux communs, latence\\
\code{core/scanorder.py} & pointages de veille d'un volume, décalages, ordre de Morton à bits renversés\\
\code{core/compute.py} & débits, faisceaux de réception, charge de traitement, tampons, puissance et latence de la chaîne
numérique\\
\code{core/repoint.py} & re-pointage en vol : erreur de pointage de la navigation, sa perte, migration entre faces\\
\code{core/allocation.py} & répartition de la puissance entre les faces par niveau de priorité\\
\code{core/margin.py} & réemploi de la marge selon le mode de réallocation, table de saturation de la poursuite\\
\code{core/redundancy.py} & redondance entre secteurs (avant le calcul, et avec les plans du calcul)\\
\code{core/validation.py} & verdict de la mission : table des contraintes, ressource limitante, état de saturation\\
\code{core/cost\_estimate.py} & estimation rapide pour l'interface\\
\code{core/mission.py} & modèle de mission, validation, JSON / CSV\\
\code{core/geometry.py} & trame de la sphère, vue 3D du vaisseau\\
\code{gui/} & interface PyQt5 avec vues 3D (PyVista)\\
\bottomrule
\end{tabularx}
\end{center}
Python 3.11 ou plus récent, NumPy, et SciPy pour le modèle de piste (équation de Riccati discrète, exponentielle de
matrice ; les fonctions de Bessel des diagrammes sont codées directement) ; Numba, s'il est installé, compile les sommes
du bruit des corps (mêmes formules en double précision, en parallèle ; la voie NumPy donne le même résultat à
$10^{-14}$ près). L'optimiseur exécute les secteurs, et les fréquences d'un secteur avec tous leurs pas de grille, dans
des processus de travail ; les lois de détection, les lois de perte de position et les inversions de $\Pd$ cumulée sont
mises en cache exactement.

\subsection{Tests}
Le banc de tests a deux familles : les sorties figées de chaque mission de test comparées à une tolérance relative de
$10^{-9}$ (non-régression), et des lois physiques et invariants qui doivent tenir quelles que soient les valeurs
($R\propto\sigma^{1/4}$, exactitude de chaque raccourci, cohérence des caches, lois de détection contre les formes
closes, \dots). Chaque regel des sorties figées est consigné avec sa cause. Un rapport d'impact donne les chiffres clés
et les formes d'onde qui changent avant tout regel. Au modèle V6.27.0, le banc compte 250 tests (dont 12 comparent les
sorties figées des cas de test), exécutés à trois niveaux : tests rapides, tests touchés par un changement, et passe
complète (table de coût, chaque cas de test calculé une fois et gardé avec une empreinte du code, rapport d'impact, banc
entier ; environ 35 minutes depuis la passe P7 de la 5.2 (table de saturation), avec une barre de progression globale).

%=====================================================================================================================
\section{Feuille de route et points ouverts}\label{sec:roadmap}
%=====================================================================================================================
\begin{center}
\begin{tabularx}{\linewidth}{@{}lL@{}}
\toprule
Phase & Étapes (état au 27 octobre 2026)\\
\midrule
0, fondations & fichier de paramètres unique, banc de tests, revue de chaque valeur --- fait\\
1, antenne et RF & diagramme réel ; désadaptation active, lobes de réseau et de jeu ; TRM et chaîne d'émission ; bande ;
chaîne de réception, bruit du ciel, convertisseurs, isolation --- fait\\
2, signal et détection & formes d'onde paramétriques ; traitement ; critère de détection ; modèles de cibles ;
environnement --- fait\\
3, fonctions radar & 3.0 repères et navigation ; 3.1 veille S ; 3.2 confirmation Ku ; 3.3 poursuite ; 3.4 erreurs du
réseau et étalonnage --- fait (modes dégradés et pannes de TRM reportés en fin de développement)\\
4, ressources & 4.1 puissance : budget partagé du vaisseau, consommation permanente, alimentation des TRM, stockage de
face --- fait ; 4.2 thermique : fluide, plaque froide, jonction, rejet de chaleur du vaisseau --- fait ; 4.3
ordonnanceur : priorités, files, chronologie synthétique, créneaux communs, latence --- fait ; 4.3 bis ordre des
pointages dans un volume --- fait ; 4.4 calcul et données --- fait ; 4.5 re-pointage en vol --- fait\\
5, mission et validation & 5.1 modèle de mission enrichi --- fait ; 5.2 optimiseur global (répartition de la
puissance, couplages, boucle globale, réemploi de la marge, tables pour le vol, réallocation et saturation) ---
fait ; 5.3 moteur de
validation (verdict et conseil faits) ; simulation Monte-Carlo ; interface opérateur ; vérification et documentation ; fin de développement\\
\bottomrule
\end{tabularx}
\end{center}

Principales dispositions provisoires, chacune affectée à une étape ultérieure :
\begin{itemize}
  \item \textbf{Ordonnanceur} : multiplexage en temps et superposition avec un entrelacement imparfait ; la part P5
        d'une veille comme classe à part ; GO-CFAR près des limbes compté (seulement conseillé aujourd'hui) ; décalage
        de la grille de veille à chaque trame (une règle) ; alignement S/Ku sur des références externes ; percée Ku
        autorisée en vol (autorisée par un interrupteur de mission depuis l'étape 5.1).
  \item \textbf{Thermique} : rejet de chaleur du vaisseau réservé au radar (\SI{6}{MW} provisoires).
  \item \textbf{Calcul} : les bilans qui dimensionnent le plan de veille gardent le plafond de \SI{50}{kW} du traitement
        (la répartition de la puissance utilise la consommation réelle) ; coût d'une resélection de forme d'onde en vol
        (valeur provisoire).
  \item \textbf{Re-pointage en vol} : vitesse et accélération maximales par mission (réglages du système aujourd'hui) ;
        enveloppe complète des attitudes (5.3) ; vaisseau en rotation dans la vue 3D (phase 6).
  \item \textbf{Mission et validation (5.x)} : les tables de présentation et les dispersions sont des valeurs
        provisoires (D) ; tables pour le vol Ku ; autres usages de la marge (veille plus loin, volumes de moindre
        priorité) ; zones de recouvrement entre faces ; recalage de l'estimation rapide sur le modèle de piste ;
        l'interface de préparation de mission est provisoire (une IHM complète vient avec la simulation globale).
  \item \textbf{Antenne} : diagramme élémentaire et couplage mutuel aux bords des tuiles (éléments isotropes) ; les
        couplages entre faces sont des ordres de grandeur optimistes à confirmer par une simulation électromagnétique.
  \item \textbf{Fin de développement} : modes dégradés (face, TRM, puissance ou ouverture Ku perdus) et pannes de TRM
        au fil d'une mission ; couplages thermiques (pompes dérivées du débit, dérive d'étalonnage liée à la
        température, dilatation des jeux de \SI{0.5}{mm}).
\end{itemize}

\appendix
%=====================================================================================================================
\section{Principales valeurs dérivées}
%=====================================================================================================================
\begin{center}
\begin{tabularx}{\linewidth}{@{}lL@{}}
\toprule
Grandeur & Valeur\\
\midrule
Positions de veille par face à \SI{2.8}{GHz} & environ \num{1664} sur le secteur complet (grille de 1,4 largeur de
faisceau d'émission)\\
Faisceaux de réception par position & 6,1 / 5,3 / 4,8 à \SIlist{-25;-30;-35}{dB} (écart de 0,5 largeur de faisceau de
réception)\\
Perte moyenne de grille d'émission sur une maille & \SI{3.4}{dB} (diagramme réel) ; grille de réception \SI{0.40}{dB}\\
Puissance d'émission disponible par face (référence \SI{2}{MW}) & \SI{1.784}{MW} électriques (chaîne 0,945, permanent
\SI{106.9}{kW})\\
Puissance électrique crête d'une visée à pleine puissance & \SI{12.3}{MW} S, \SI{14.1}{MW} Ku (stockage de face
\SI{5}{MJ})\\
Impulsion la plus longue & \SI{2.0}{ms} en S (\SI{300}{km} aveugles), \SI{2.08}{ms} en Ku\\
PAE à $d=0,65$, impulsion longue & 0,778 en S, 0,621 en Ku\\
Facteur de bruit du CAN & \SI{20.0}{dB} (13 bits, \SI{+10}{dBm}, \SI{250}{MS/s})\\
Bruit au CAN & \SI{-56.4}{dBFS} en S\\
Portée minimale sans saturation (S, dans l'axe) & cible de référence \SI{12.2}{km}, KKV \SI{2.3}{km}, vaisseau
\SI{9.1}{km}\\
Pondération en distance \SIlist{-25;-30;-35}{dB} & perte \SIlist{0.41;0.68;0.93}{dB}, élargissement $\times$1,18 / 1,27 /
1,34\\
Doppler à \SI{25}{km/s} & \SI{470}{kHz} en S ; NLFM possible au-delà de \SI{23}{MHz} en S, \SI{117}{MHz} en Ku\\
Désadaptation active aller-retour à \SI{2.8}{GHz} & \SI{0.26}{dB} à \SI{0}{\degree}, 0 à \SI{45}{\degree},
\SI{1.12}{dB} à \SI{70}{\degree}, \SI{2.11}{dB} à \SI{75}{\degree}\\
Aire du secteur en $(u,v)$ & 2,0006\\
\bottomrule
\end{tabularx}
\end{center}

%=====================================================================================================================
\section{Glossaire : les grandeurs et la physique qu'elles traduisent}\label{sec:glossary}
%=====================================================================================================================
Chaque grandeur principale du modèle, son sens et le phénomène physique qu'elle traduit. Les valeurs sont celles du
fichier de paramètres à la date du document.

\subsection{Géométrie, antenne et faisceaux}
{\small
\begin{longtable}{@{}p{0.2\linewidth}p{0.76\linewidth}@{}}
\toprule
Grandeur & Sens et physique\\
\midrule
\endhead
$\lambda$, $f$ & Longueur d'onde et fréquence porteuse, $\lambda = c/f$ (\SI{10.7}{cm} à \SI{2.8}{GHz}, \SI{2.14}{cm} à
\SI{14}{GHz}). Toute taille angulaire du réseau varie en $\lambda/D$.\\
$\Denv$, $x_{\mathrm{rms}}$ & Enveloppe de l'ouverture (\SI{4.17}{m}) et étendue RMS des positions des éléments selon un
axe. $\lambda/\Denv$ fixe la largeur du faisceau ; $x_{\mathrm{rms}}$ fixe la meilleure précision angulaire possible
(une ouverture mesure la pente de phase du front d'onde sur son étendue).\\
$(u, v)$, $\theta$ & Cosinus directeurs dans le plan de la face et angle de dépointage par rapport à la normale. Un
réseau plan forme un faisceau de largeur constante en $(u,v)$ ; en angle, il s'élargit en $1/\cos\theta$ parce que
l'ouverture projetée diminue.\\
Largeur à mi-puissance & Angle entre les deux directions où le gain est divisé par deux. Émission (éclairement
uniforme) $1,036\,\lambda/\Denv$ ; réception plus large à cause de la pondération.\\
$G$, $\eta$ & Gain $4\pi A\eta\cos\theta/\lambda^2$ : combien le réseau concentre la puissance dans la direction du
faisceau par rapport à un rayonneur isotrope. $\eta$ est le rendement d'éclairement : 1 pour une excitation uniforme,
moins avec une pondération qui abaisse les lobes secondaires (énergie retirée des bords).\\
Niveau des lobes secondaires & Gain relatif hors du lobe principal. Des lobes bas rejettent les sources fortes hors de
l'axe (Soleil, Terre, brouilleurs, fouillis) au prix du gain. La pondération de réception (loi de Taylor) le fixe par
volume.\\
Lobe de réseau & Réplique complète du lobe principal qui apparaît quand l'écart entre éléments dépasse environ
$\lambda/(1+\sin\theta)$ : il fixe le haut de la bande utilisable et la limite de dépointage.\\
Désadaptation active $\Gamma$ & Dans un réseau dépointé, les éléments se couplent entre eux ; l'impédance vue par chaque
TRM change avec le dépointage, une partie de la puissance est réfléchie. La couche WAIM adapte le réseau à
\SI{45}{\degree} ; la perte monte à l'approche du lobe de réseau.\\
TRSA, jeux & Tuile de sous-réseau d'émission-réception et bandes mortes entre tuiles. La rupture du pas régulier aux
jeux crée des « lobes de jeu » que la pondération ne peut abaisser.\\
Pas de grille & Distance entre faisceaux de veille voisins, en largeurs de faisceau d'émission (1,3 à 2,0). Un pas plus
large donne moins de positions mais une cible plus loin du centre du faisceau le plus proche (perte de grille).\\
\bottomrule
\end{longtable}}

\subsection{Émetteur et puissance}
{\small
\begin{longtable}{@{}p{0.2\linewidth}p{0.76\linewidth}@{}}
\toprule
Grandeur & Sens et physique\\
\midrule
\endhead
$P_{\mathrm{pk}}$ & Puissance RF crête d'un TRM à la sortie du PA (\SI{2}{kW} S, \SI{75}{W} Ku). Les PA fonctionnent
saturés (enveloppe constante), le régime le plus efficace.\\
PAE $\eta$ & Rendement en puissance ajoutée : puissance RF sortante sur puissance continue entrante. Abaissé par les
fronts d'impulsion (courant continu consommé pendant l'établissement de la RF) et par l'échauffement du canal à fort
facteur de charge.\\
$\tau$, PRI, $d$ & Durée d'impulsion, intervalle de répétition et facteur de charge $d = \tau/\mathrm{PRI}$. Les
impulsions longues portent plus d'énergie (portée) mais aveuglent le récepteur plus longtemps (distance aveugle
$c\tau/2$).\\
$N$, train & Nombre d'impulsions d'une visée ; leurs échos sont intégrés de façon cohérente (phase gardée) ou non
cohérente (puissances additionnées).\\
Réserve d'énergie & Batterie de condensateurs locale de chaque TRM qui délivre l'impulsion ; la tension baisse à mesure
qu'elle se vide (affaissement) et la sortie du PA suit.\\
Alimentation des TRM & Puissance continue qui recharge la réserve entre les impulsions ; elle borne la puissance moyenne
d'un train.\\
Rendement de chaîne & Rendement de l'alimentation et de la distribution (0,945) : la part de la puissance du bus qui
atteint les TRM et les récepteurs.\\
Puissance permanente & Récepteurs, convertisseurs, formation de faisceaux, contrôle : consommée que la face émette ou
non.\\
Stockage de face & Réserve d'énergie qui délivre la puissance crête d'une visée (tous les TRM à la fois) au-dessus de ce
que le bus fournit en moyenne.\\
Température de jonction & Température du canal du transistor d'un PA ; elle monte à chaque impulsion (pôle rapide) et
sur le train (pôle lent) ; au-delà d'environ \SI{225}{\celsius}, la durée de vie du GaN chute.\\
Plaque froide, flux & Plaque à microcanaux où le fluide prend la chaleur ; le flux qu'elle accepte par unité de surface
vaut $h \times$ facteur de surface mouillée $\times$ écart de température admis.\\
Inertie thermique & Chaleur que le fluide et la structure stockent par kelvin : elle absorbe les rafales d'émission.\\
Conversion numérique vers le bas & Ramener les échantillons d'une voie à la bande du signal, près du convertisseur,
pour réduire le débit.\\
Formation de faisceaux par sous-réseau & Faisceaux partiels formés dans une TRSA et sommés par le processeur de face.\\
Présentation & Direction d'où le radar voit la cible, dans les axes propres de la cible (azimut et élévation depuis son
nez).\\
Message de navigation & Attitude, vitesse et orbite du vaisseau envoyées au calculateur radar (ici 100 par seconde).\\
Loi de Little & Nombre moyen d'éléments en cours = cadence d'arrivée $\times$ durée de séjour de chacun ; ici, les
fenêtres d'écoute ouvertes.\\
Ressource dimensionnante & Ressource la plus chargée d'une face (énergie, temps, chaleur, calcul, faisceaux, mémoire) ;
le réemploi de la marge la porte jusqu'au plafond \code{load\_ceiling}.\\
Saturation gérable, critique & Pistes au-delà de la spécification tenues sans dégrader la veille ; au-delà, dégradation
puis abandon de tâches selon le mode de réallocation.\\
\bottomrule
\end{longtable}}

\subsection{Récepteur, bruit et interférences}
{\small
\begin{longtable}{@{}p{0.2\linewidth}p{0.76\linewidth}@{}}
\toprule
Grandeur & Sens et physique\\
\midrule
\endhead
$\Tsys$ & Température de bruit du système : la puissance de bruit thermique $kT_{\mathrm{sys}}B$ qu'ajoute le récepteur,
ramenée à l'antenne. Elle somme le ciel vu par l'antenne, les pertes avant le premier amplificateur (un élément à pertes
à la température $T$ rayonne $(L-1)T$), le bruit de l'amplificateur et celui du convertisseur, chacun divisé par le
gain qui le précède (Friis).\\
Facteur de bruit, ENOB & Facteur de bruit : combien un étage dégrade le SNR. ENOB : nombre effectif de bits du
convertisseur ; son bruit de quantification donne le facteur de bruit du convertisseur.\\
Écrêtage, dynamique & Un convertisseur sature au-dessus de sa pleine échelle ; chaque élément étant numérisé avant la
formation des faisceaux, un écho fort ou une surface brillante dans le lobe principal peut écrêter tous les
convertisseurs à la fois.\\
Bruit du ciel $T_b$ & Tout corps émet un bruit thermique à sa température de brillance (Soleil \SI{1e5}{K} en S, Terre
\SI{290}{K}, Lune \SI{230}{K}, galaxie quelques K). Vu à travers le diagramme, il élève $\Tsys$ ; un corps dans le lobe
principal peut aveugler une direction.\\
Polarisation, $\rho$ & La polarisation de l'écho est inconnue ; deux voies de réception en recueillent les deux
composantes. $\rho$ est le rapport de SER croisée sur principale de la cible : l'énergie supplémentaire que recueille la
seconde voie.\\
Rotation de Faraday & Les électrons de l'ionosphère dans le champ magnétique tournent le plan de polarisation en
$\propto B_\parallel\,\mathrm{TEC}/f^2$ ; sans effet sur la détection avec deux voies, une source d'erreur pour la
polarimétrie.\\
Fouillis & Échos d'une surface (Terre, Lune) à la même distance que la cible, par les lobes secondaires, en concurrence
avec elle. Le filtrage Doppler le rejette quand les vitesses diffèrent.\\
Brouilleur & Bruit envoyé vers le radar : à distance (dans les lobes secondaires, annulé adaptativement) ou
d'autoprotection (porté par la cible, dans le lobe principal). Portée de percée : là où l'écho domine de nouveau le
brouillage.\\
\bottomrule
\end{longtable}}

\subsection{Forme d'onde, traitement et détection}
{\small
\begin{longtable}{@{}p{0.2\linewidth}p{0.76\linewidth}@{}}
\toprule
Grandeur & Sens et physique\\
\midrule
\endhead
$B$, compression d'impulsion & Largeur de bande du signal. Une impulsion longue codée, comprimée à la réception, a
l'énergie de l'impulsion longue et la résolution en distance $c/2B$ d'une courte. LFM : balayage linéaire de
fréquence, tolérant au Doppler ; NLFM : balayage non linéaire dont la forme de spectre donne des lobes secondaires en
distance bas sans perte de pondération.\\
Lobes secondaires en distance, pondération & Réponse résiduelle de l'impulsion comprimée à côté de son pic ; une
pondération les abaisse (pour garder visible une cible faible à côté d'une forte) au prix du SNR et de la résolution.\\
Doppler $f_d$ & Décalage de fréquence $2v/\lambda$ de l'écho d'une cible qui se rapproche à $v$ (\SI{470}{kHz} à
\SI{25}{km/s} en S). Il décale le pic comprimé (couplage) et sépare les cibles du fouillis.\\
Bloc cohérent, $T_0$ & Partie d'un train intégrée avec ses phases. Cela fonctionne tant que la phase de l'écho reste
prévisible : une cible en rotation se décorrèle en son temps de cohérence $T_0 = \lambda/(2L\omega)$, et le bruit de
phase d'un oscillateur la dégrade aussi.\\
Migration en distance & Une cible rapide traverse plusieurs cases distance pendant la visée ; elle est compensée, au
prix d'un banc d'hypothèses de vitesse.\\
Bancs de vitesses, d'accélérations & Le traitement essaie plusieurs hypothèses de vitesse (et d'accélération) ; chacune
coûte un filtre (calcul) et compte comme des cases supplémentaires pour les fausses alarmes ; une cible entre deux
hypothèses perd du SNR.\\
Chevauchement & Une cible qui tombe entre deux échantillons de la grille distance ou Doppler perd du SNR ; le
suréchantillonnage le réduit.\\
CFAR & Détecteur à taux de fausses alarmes constant : le seuil est fixé à partir du bruit mesuré dans les cases
voisines ; l'aléa de cette estimation coûte une petite perte.\\
$\Pd$, $\Pfa$ & Probabilité de détecter une cible et de fausse alarme par case indépendante.\\
Swerling, forme Gamma $k$ & Modèle statistique de la puissance de l'écho d'une visée à l'autre : stable (sphère),
Swerling 1/2 (nombreux réflecteurs, exponentielle), Swerling 3/4 (un réflecteur dominant). Le modèle utilise une loi
Gamma de forme $k$ : plus $k$ est grand, plus l'écho est stable ; les échantillons indépendants (impulsions,
fréquences, polarisations) augmentent $k$.\\
$\SNR_{\mathrm{req}}$ & SNR dont le détecteur a besoin pour atteindre la $\Pd$ exigée à la $\Pfa$ et la fluctuation
données.\\
$\Pd$ cumulée, $q$ & Probabilité qu'une cible qui se rapproche ait été détectée au moins une fois quand elle atteint la
portée exigée ; $q = vRR/R$ est la distance relative qu'elle parcourt entre deux visées.\\
Perte de positions & SNR perdu parce que la cible est quelque part dans la case de veille (centre du faisceau, grilles
distance et Doppler, filtre de vitesse), moyenné en $\Pd$.\\
\bottomrule
\end{longtable}}

\subsection{Confirmation et poursuite}
{\small
\begin{longtable}{@{}p{0.2\linewidth}p{0.76\linewidth}@{}}
\toprule
Grandeur & Sens et physique\\
\midrule
\endhead
Confirmation & Une visée Ku pointée sur une détection S : le faisceau Ku étroit (5 fois plus étroit) vérifie qu'elle est
réelle avant l'ouverture d'une piste ; l'erreur angulaire de l'estimation S décide s'il faut un faisceau ou une grappe
de 7.\\
Filtre de Kalman, modèle de Singer & Estimateur récursif de la position, de la vitesse et de l'accélération. Le modèle
de Singer traite l'accélération de la cible comme un processus aléatoire d'écart-type $\sigma_m$ et de constante de
temps $\tau_m$ (la durée d'une man\oe{}uvre).\\
Erreur prédite, filtrée & Incertitude sur la cible juste avant une mise à jour (où pointer) et juste après (qualité de
piste délivrée).\\
Perte de pointage & Le faisceau de mise à jour est pointé sur la prédiction ; la cible est hors du centre du faisceau de
l'erreur de prédiction et perd du gain en émission et en réception.\\
Extrapolation, perte de piste & Une mise à jour manquée (pas de détection) laisse la piste extrapoler une période de
plus ; l'erreur grandit, la visée suivante est plus difficile ; quelques manques d'affilée perdent la piste.\\
Discrimination & Distinguer une charge (KKV) d'un leurre : polarisation (taux de dépolarisation), statistique de
fluctuation, niveau de SER (optique d'autodirecteur) et man\oe{}uvres diffèrent.\\
Erreurs du réseau, étalonnage & Erreurs résiduelles d'amplitude et de phase des TRM après étalonnage : elles baissent le
gain (Ruze) et élèvent un plancher de lobes secondaires aléatoires ; elles dérivent avec la température entre deux
étalonnages.\\
\bottomrule
\end{longtable}}

%=====================================================================================================================
\section{Effets modélisés et leur impact}\label{sec:effects}
%=====================================================================================================================
Chaque effet pris en compte par le modèle, ce qu'il représente, et la plage de son impact. Sauf mention, l'impact est
une perte de SNR en dB (un gain si négatif) : une perte de $x\dB$ raccourcit la portée du facteur $10^{-x/40}$
(\SI{1}{dB} $\approx -5,6\,\%$, \SI{3}{dB} $\approx -16\,\%$) ou, à portée fixée, multiplie l'énergie par $10^{x/10}$.
Les plages viennent du fichier de paramètres et d'un échantillon de l'espace de recherche S (3 secteurs de test, 900
formes d'onde, 3 fréquences, dépointage de 0 à \SI{70}{\degree}) et de l'espace des visées Ku (300 visées).

{\small
\begin{longtable}{@{}p{0.22\linewidth}p{0.47\linewidth}p{0.25\linewidth}@{}}
\toprule
Effet & Ce qu'il représente & Impact (plage)\\
\midrule
\endhead
\multicolumn{3}{@{}l}{\textbf{Antenne et chaîne d'émission}}\\
Ouverture projetée $\cos\theta$ & Le réseau vu de la direction de dépointage est plus petit : gain en émission et en
réception & 0 à \SI{9.3}{dB} aller-retour à \SI{70}{\degree}\\
Désadaptation active & Réflexion de la puissance des TRM par l'impédance couplée du réseau, fonction du dépointage,
WAIM adaptée à \SI{45}{\degree}, montant près des lobes de réseau & 0 (\SI{45}{\degree}) à \SI{1.1}{dB} à
\SI{70}{\degree}, \SI{2.1}{dB} à \SI{75}{\degree} (\SI{2.8}{GHz}, aller-retour)\\
Chaîne d'émission & Filtre d'harmoniques, circulateur, coupleur, ligne entre PA et élément & \SI{0.60}{dB} S,
\SI{0.95}{dB} Ku\\
Occultation par l'autre bande & Les éléments de l'autre bande diffusent une partie du champ & \SI{0.2}{dB} S,
\SI{0.3}{dB} Ku\\
Erreurs du réseau (étalonnage) & Erreurs résiduelles d'amplitude et de phase des TRM et dérive thermique entre
étalonnages (Ruze) & \SI{0.006}{dB} ; plancher de lobes \SI{-59}{dB}\\
Option de pondération de réception & Pondération de Taylor à \SIlist{-25;-30;-35}{dB} de lobes secondaires & gain 0,
$-0,53$, \SI{-1.00}{dB}\\
Bords de bande & Dégradation en cosinus surélevé du TRM hors de la bande c\oe{}ur, moyennée sur le spectre occupé & 0 (c\oe{}ur)
à \SI{1.0}{dB} S / \SI{0.5}{dB} Ku aux bords (échantillon : 0--0,17)\\
Affaissement de la réserve & Sortie du PA qui suit la tension de la réserve pendant l'impulsion & 0 à \SI{0.46}{dB}\\
PAE (fronts, facteur de charge) & Puissance continue consommée sans RF pleine sur les fronts, échauffement du canal à
fort facteur de charge & énergie $+0,1$ à $+2,8\,\%$ S, jusqu'à $+4,4\,\%$ Ku\\
Multifaisceau en émission & Intermodulation de la superposition en phase seule ; partition : gain plus faible, faisceau
plus large & 0,4 / 0,6 / \SI{0.8}{dB} pour 2 / 4 / 8 faisceaux\\
\midrule
\multicolumn{3}{@{}l}{\textbf{Forme d'onde et traitement}}\\
Compression d'impulsion & Désadaptation NLFM, ou pondération LFM plus perte Doppler & \SI{0.1}{dB} (NLFM) ; LFM 0,41 à
\SI{0.93}{dB} + Doppler (échantillon Ku jusqu'à \SI{1.26}{dB})\\
Pondération Doppler & Pondération de Taylor d'un bloc cohérent contre le fouillis & \SI{0.68}{dB} (\SI{-30}{dB}),
\SI{1.14}{dB} (\SI{-40}{dB})\\
Éclipse & Échos perdus pendant que le récepteur est sourd lors des impulsions suivantes d'un train ambigu en distance &
$-10\log_{10}(1-(\tau+t_{\mathrm{rec}})/\mathrm{PRI})$, typiquement 0,1 à \SI{1}{dB}\\
Cohérence & Décorrélation de la cible et bruit de phase de l'oscillateur dans un bloc cohérent & 0 à \SI{0.63}{dB}\\
Migration en distance & Résidu de la compensation quand la cible traverse plusieurs cases & 0 ou \SI{0.1}{dB}\\
Banc de vitesses & Cible entre deux hypothèses de vitesse (moyenne sur le pas) & 0 à \SI{0.34}{dB}\\
Accélération & Phase quadratique et courbure en distance d'une cible qui accélère, ou résidu du banc & 0 à
\SI{0.35}{dB}\\
Chevauchement distance et Doppler & Cible entre deux échantillons de grille (grilles suréchantillonnées $\times4$) &
\SI{0.05}{dB} et \SI{0.07}{dB}\\
CFAR & Aléa de l'estimation du bruit (CA, 1024 cases ; GO aux bords de fouillis) & \SI{0.04}{dB} (GO
\SI{0.05}{dB})\\
Précision finie & Quantification et arithmétique du traitement numérique & \SI{0.3}{dB}\\
\midrule
\multicolumn{3}{@{}l}{\textbf{Détection}}\\
Grille de veille (émission) & Cible n'importe où dans la maille hexagonale de la grille de veille, diagramme réel &
\SI{3.3}{dB} en moyenne (pas 1,4), plus pour des pas plus larges\\
Grille des faisceaux de réception & Cible entre les faisceaux de réception numériques (écart 0,5 largeur de faisceau de
réception) & \SI{0.40}{dB}\\
Perte de positions (totale) & Tous les effets de position ci-dessus à travers la $\Pd$ moyennée sur la case & 3,1 à
\SI{4.3}{dB} (veille), 0,05 à \SI{0.5}{dB} (visées Ku)\\
Fluctuation de la cible & Puissance de l'écho variant d'une visée à l'autre ; stable, Swerling 1, Swerling 3 à
$\Pd = 0,9$, $\Pfa = 10^{-8}$, un échantillon & $\SNR_{\mathrm{req}}$ 14,2 / 22,4 / \SI{18.5}{dB} (la fluctuation
coûte jusqu'à \SI{8}{dB} ; agilité et diversité en récupèrent une partie)\\
Politique de fausses alarmes & $\Pfa$ par case tirée des fausses alarmes tolérées, $10^{-14}$ à $10^{-4}$ &
$\SNR_{\mathrm{req}}$ $\pm$\SI{3}{dB} autour de $10^{-8}$\\
Critère cumulé & Plusieurs visées d'une cible qui se rapproche au lieu d'une & $-6,8\dB$ ($q = 0,1$) à $+6,0\dB$
($q = 1$) par rapport à une visée\\
Polarisation (GLRT) & Polarisation inconnue de l'écho estimée sur la visée & 0,03 à \SI{0.47}{dB}\\
Énergie en polarisation croisée & La seconde voie recueille l'écho en polarisation croisée & gain $-0,04$ (KKV) à
$-1,76\dB$ (débris)\\
\midrule
\multicolumn{3}{@{}l}{\textbf{Récepteur et environnement}}\\
Bruit du système (ciel froid) & Cascade de Friis : ciel, pertes avant LNA, LNA, après LNA, convertisseur & $\Tsys$
\SI{144.5}{K} S, \SI{230.8}{K} Ku\\
Fond galactique & Brillance de la galaxie à travers le diagramme (pire cas sans repère) & jusqu'à $+$\SI{8}{K}
($\approx +$\SI{0.2}{dB} en S)\\
Corps (Soleil, Terre, Lune) & Température de brillance à travers le diagramme de réception & Terre $+$2,5 à
\SI{4}{dB} S ; Lune $+$\SI{2.6}{dB} Ku ; Soleil $\approx +$\SI{17}{dB} : aveuglé au-delà de \SI{10}{dB}\\
Structure du vaisseau & Occultation du tube du faisceau par la coque ; émission de la coque et des radiateurs & 0 avec
l'inclinaison de \SI{25}{\degree} ; coque $+$\SI{2.9}{K}\\
Fouillis de surface & Échos de la Terre ou de la Lune dans la case distance de la cible par les lobes secondaires, après
filtrage Doppler & 0 à l'aveuglement (augmentation de SNR $G$ par forme d'onde) ; exigence ramenée à la surface\\
Brouilleurs à distance & Bruit par les lobes secondaires, annulé adaptativement (\SI{40}{dB}) & montée de $\Tsys$,
fonction de la direction\\
Brouilleur d'autoprotection & Bruit de barrage dans le lobe principal & portée limitée à la portée de percée (visée de
percée Ku dans l'axe du brouillage)\\
Rotation de Faraday & Rotation de la polarisation par l'ionosphère & 0 pour la détection ; erreur en polarimétrie\\
Écrêtage des convertisseurs & Écho fort ou surface brillante saturant le convertisseur de chaque élément & portée
minimale sans saturation (KKV en S \SI{2.3}{km}) ; l'atténuateur après LNA l'abaisse jusqu'à $\times5,6$\\
Distance aveugle & Récepteur sourd pendant l'impulsion et la récupération du limiteur & $c(\tau+t_{\mathrm{rec}})/2$ :
\SI{3}{km} (\SI{20}{\micro s}) à \SI{300}{km} (\SI{2}{ms})\\
\midrule
\multicolumn{3}{@{}l}{\textbf{Confirmation et poursuite}}\\
Erreur angulaire S vers Ku & Erreur de l'estimation S et biais S/Ku face au faisceau Ku étroit & marges de confirmation
\SIrange{0}{12}{dB}, visée seule ou grappe de 7\\
Perte de pointage de piste & Erreur de prédiction de la piste face au faisceau Ku (émission et réception) & 0,4 (0,1
largeur de faisceau) à \SI{3.9}{dB} (0,3 largeur) ; maille de grappe \SI{3.8}{dB} sur 7 visées\\
Bancs de poursuite & Hypothèses de vitesse et d'accélération sur $\pm3\sigma$ de la prédiction & 0 à
$\approx$\SI{0.4}{dB}\\
Couplage distance--Doppler & Résidu du couplage après correction Doppler, sur la précision en distance & élève le SNR
des mises à jour pilotées par la distance (impulsions longues, bande étroite)\\
Perte et initialisation de piste & Mises à jour manquées, extrapolation, réinitialisation ; visées d'initialisation
après confirmation & perte $10^{-3}$ à $10^{-2}$ par mise à jour ; initialisation 1 à $\approx 30$ visées\\
Planchers de mesure & Étalonnage, erreurs des TRM, horloge, oscillateur & \SI{0.05}{mrad}, \SI{0.05}{m},
\SI{0.01}{m/s}\\
Discrimination & Quatre indices sur les visées de poursuite & visées nécessaires pour atteindre $P = 0,95$ (optimiste
sans dispersion de classe)\\
\midrule
\multicolumn{3}{@{}l}{\textbf{Ressources}}\\
Temps d'étalonnage & Étalonnage en vol pris sur le temps d'émission & \num{3.3e-5} (S), \num{1.7e-4} (Ku) du temps\\
Puissance permanente & Convertisseurs, chaînes de réception, formation de faisceaux, contrôle & \SI{106.9}{kW} par face
($\approx 5\,\%$ de \SI{2}{MW}) ; pompes 2\,\% de la chaleur\\
Alimentation des TRM & Puissance moyenne et recharge de la réserve sur un train & retire des formes d'onde (0 en S,
3\,\% des visées Ku)\\
Stockage de face & Puissance crête d'une visée au-dessus de la moyenne du bus & règle d'ordonnanceur
$\approx$\SI{0.3}{s} d'émission contiguë\\
Créneaux d'émission communs & Toutes les faces émettent une bande ensemble ; récepteurs sourds dans les créneaux &
éclipse $-10\log_{10}(1-\rho)$ : \SI{0.1}{dB} sur une mission de test (\SI{2.4}{\percent} du temps)\\
Files à priorités & Attente d'une tâche derrière les priorités plus hautes, sans préemption & latence de confirmation
\SIrange{8}{80}{ms} (95\,\%) sur une mission de test\\
Charge de traitement & Formation de faisceaux, compression, bancs, CFAR, annulation contre la capacité de la puissance
plafond & $\le 0,6\,\%$ de la capacité sur les missions de test : non limitant\\
Puissance du traitement & Part statique plus énergie par opération et par bit & environ \SI{10}{kW} par face au lieu du
plafond de \SI{50}{kW} (donnée en sortie ; le dimensionnement garde le plafond)\\
Latence de traitement & Flux de données et traitement après un passage & \SI{0.01}{\micro s} à \SI{5}{ms}, dans la
latence de confirmation\\
Pointage de navigation & Connaissance d'attitude et de vitesse, accélération non mesurée entre messages &
\num{3e-5}\,dB en S, \num{8e-4}\,dB en Ku ($\le 0,006$ largeur de faisceau)\\
Présentation de la cible & SER en fonction de l'azimut et de l'élévation de la présentation, la pire (KKV de face) &
aucune portée de détection changée aujourd'hui ; indice de SER de discrimination $\div 3$ (KKV / leurre, 5 visées)\\
Refroidissement d'une face & Plaque froide à microcanaux, fluide diélectrique & capacité \SI{15.9}{MW} par face : non
limitant (dissipation $\le$ \SI{2}{MW})\\
Température de jonction & Échauffement des transistors des PA par les impulsions & $\le$\SI{63}{\celsius} contre
\SI{225}{\celsius} : non limitant\\
Rejet de chaleur du vaisseau & Chaleur du radar face au rejet réservé & par exemple \SI{1.0}{MW} sur \SI{6}{MW}
(mission de test)\\
Réemploi de la marge & Revisites (\code{PRIO\_SEARCH}) ou cadence de poursuite (\code{PRIO\_TRACK}) jusqu'au plafond de la ressource dimensionnante, plan gardé
& portée cumulée $+19$ à $+89\,\%$ sur une mission de test ($X = \SI{80}{\percent}$)\\
Réglages de calcul & Approximations de l'optimiseur (lois arrondies, peu de fréquences dans un spectre calme, \dots) &
$\le$\SI{0.003}{dB} sur $\SNR_{\mathrm{req}}$ ; énergie de veille $-5$ à $+9\,\%$ par secteur\\
\bottomrule
\end{longtable}}

\end{document}
